\begin{filecontents*}[overwrite]{\jobname.ist}
headings_flag 1
numhead_positive "数字"
symhead_positive "符号"
heading_prefix "{\\bfseries "
heading_suffix "}\\nopagebreak\\par\\vskip 3pt "
delim_0 "，"
delim_1 "，"
delim_2 "，"
delim_n ", "
\end{filecontents*}



\usepackage{amsmath,amssymb,amsthm}
\usepackage{mathtools}
\usepackage{mathrsfs}          %
\usepackage{tikz}              %
\usepackage{tikz-cd}           %
\usepackage{etoolbox}          %
\usepackage[a4paper,margin=2.8cm]{geometry}
\usepackage{enumitem}
\setlist{itemsep=2pt,topsep=4pt}
\usepackage[font=small,labelsep=quad]{caption}  %
\emergencystretch=1.5em %
\vfuzz=3pt              %
\hbadness=2500          %
\xeCJKsetup{PunctStyle=quanjiao}%
\usepackage{array}
\usepackage{longtable,multicol,booktabs}

\ctexset{
  chapter = {
    name       = {第,章},
    number     = {\arabic{chapter}},
    format     = {\huge\bfseries\raggedright},
    nameformat = {\large\bfseries},
    numberformat={\large\bfseries},
    aftername  = {\par\vspace{14pt}},
    beforeskip = 24pt,
    afterskip  = 40pt,
  },
  section = {
    format   = {\Large\bfseries},
    aftername= {\quad},
  },
  subsection = {
    format   = {\bfseries},
    aftername= {\quad},
  },
}
\renewcommand{\thesubsection}{\thesection.\alph{subsection}}

\makeatletter
\patchcmd{\@oddhead}{\slshape}{}{}{}
\patchcmd{\@evenhead}{\slshape}{}{}{}
\makeatother

\newif\ifzhnames
\zhnamestrue
\newcommand{\trn}[2]{\ifzhnames #2\else #1\fi}

\newcommand{\trhTheorem}{定理}%
\newcommand{\trhProposition}{命题}%
\newcommand{\trhLemma}{引理}%
\newcommand{\trhCorollary}{推论}%
\newcommand{\trhConjecture}{猜想}%
\newcommand{\trhAssumption}{假设}%
\newcommand{\trhDefinition}{定义}%
\newcommand{\trhExample}{例}%
\newcommand{\trhRemark}{注记}%
\newcommand{\trhNotation}{记号}%
\newcommand{\trhWarning}{警告}%
\newcommand{\trhExercise}{练习}%
\newcommand{\trMorse}{\trn{Morse}{莫尔斯}}%
\newcommand{\trMilnor}{\trn{Milnor}{米尔诺}}%
\newcommand{\trSmale}{\trn{Smale}{斯梅尔}}%
\newcommand{\trWitten}{\trn{Witten}{维滕}}%
\newcommand{\trGromov}{\trn{Gromov}{格罗莫夫}}%
\newcommand{\trArnold}{\trn{Arnold}{阿诺德}}%
\newcommand{\trMaslov}{\trn{Maslov}{马斯洛夫}}%
\newcommand{\trConley}{\trn{Conley}{康利}}%
\newcommand{\trHamilton}{\trn{Hamilton}{哈密顿}}%
\newcommand{\trLagrange}{\trn{Lagrange}{拉格朗日}}%
\newcommand{\trLegendre}{\trn{Legendre}{勒让德}}%
\newcommand{\trPoisson}{\trn{Poisson}{泊松}}%
\newcommand{\trFourier}{\trn{Fourier}{傅里叶}}%
\newcommand{\trParseval}{\trn{Parseval}{帕塞瓦尔}}%
\newcommand{\trLiouville}{\trn{Liouville}{刘维尔}}%
\newcommand{\trLipschitz}{\trn{Lipschitz}{利普希茨}}%
\newcommand{\trFloquet}{\trn{Floquet}{弗洛凯}}%
\newcommand{\trDarboux}{\trn{Darboux}{达布}}%
\newcommand{\trCauchy}{\trn{Cauchy}{柯西}}%
\newcommand{\trRiemann}{\trn{Riemann}{黎曼}}%
\newcommand{\trEuler}{\trn{Euler}{欧拉}}%
\newcommand{\trPoincare}{\trn{Poincar\'e}{庞加莱}}%
\newcommand{\trWhitney}{\trn{Whitney}{惠特尼}}%
\newcommand{\trThom}{\trn{Thom}{托姆}}%
\newcommand{\trPontryagin}{\trn{Pontryagin}{庞特里亚金}}%
\newcommand{\trChern}{\trn{Chern}{陈}}%
\newcommand{\trHopf}{\trn{Hopf}{霍普夫}}%
\newcommand{\trBanach}{\trn{Banach}{巴拿赫}}%
\newcommand{\trHilbert}{\trn{Hilbert}{希尔伯特}}%
\newcommand{\trFredholm}{\trn{Fredholm}{弗雷德霍姆}}%
\newcommand{\trSobolev}{\trn{Sobolev}{索伯列夫}}%
\newcommand{\trArzela}{\trn{Arzel\`a}{阿尔泽拉}}%
\newcommand{\trAscoli}{\trn{Ascoli}{阿斯科利}}%
\newcommand{\trStokes}{\trn{Stokes}{斯托克斯}}%
\newcommand{\trRham}{\trn{de Rham}{德拉姆}}%
\newcommand{\trHodge}{\trn{Hodge}{霍奇}}%
\newcommand{\trGrassmann}{\trn{Grassmann}{格拉斯曼}}%
\newcommand{\trRouche}{\trn{Rouch\'e}{鲁歇}}%
\newcommand{\trHoelder}{\trn{H\"older}{赫尔德}}%
\newcommand{\trRiesz}{\trn{Riesz}{里斯}}%
\newcommand{\trYoung}{\trn{Young}{杨}}%
\newcommand{\trRellich}{\trn{Rellich}{雷利希}}%
\newcommand{\trCalderonZygmund}{\trn{Calder\'on--Zygmund}{卡尔德龙--齐格蒙德}}%
\newcommand{\trHilleYosida}{\trn{Hille--Yosida}{希勒--吉田}}%
\newcommand{\trHahnBanach}{\trn{Hahn--Banach}{哈恩--巴拿赫}}%
\newcommand{\trJacobi}{\trn{Jacobi}{雅可比}}%
\newcommand{\trLefschetz}{\trn{Lefschetz}{莱夫谢茨}}%
\newcommand{\trBorel}{\trn{Borel}{博雷尔}}%
\newcommand{\trBrouwer}{\trn{Brouwer}{布劳威尔}}%
\newcommand{\trWeyl}{\trn{Weyl}{外尔}}%
\newcommand{\trSchauder}{\trn{Schauder}{绍德尔}}%
\newcommand{\trLeray}{\trn{Leray}{勒雷}}%
\newcommand{\trLeviCivita}{\trn{Levi-Civita}{列维--奇维塔}}%
\newcommand{\trLaurent}{\trn{Laurent}{洛朗}}%
\newcommand{\trHeinz}{\trn{Heinz}{海因茨}}%
\newcommand{\trPicard}{\trn{Picard}{皮卡}}%
\newcommand{\trNewton}{\trn{Newton}{牛顿}}%
\newcommand{\trRolle}{\trn{Rolle}{罗尔}}%
\newcommand{\trSerre}{\trn{Serre}{塞尔}}%
\newcommand{\trBott}{\trn{Bott}{博特}}%
\newcommand{\trMoser}{\trn{Moser}{莫泽}}%
\newcommand{\trSard}{\trn{Sard}{萨德}}%
\newcommand{\trDonaldson}{\trn{Donaldson}{唐纳森}}%
\newcommand{\trWeierstrass}{\trn{Weierstrass}{魏尔斯特拉斯}}%
\newcommand{\trTaylor}{\trn{Taylor}{泰勒}}%
\newcommand{\trHadamard}{\trn{Hadamard}{阿达马}}%
\newcommand{\trAudin}{\trn{Audin}{奥丹}}%
\newcommand{\trDamian}{\trn{Damian}{达米安}}%
\newcommand{\trFloer}{\trn{Floer}{弗洛尔}}%
\newcommand{\trAudinFull}{\trn{Mich\`ele Audin}{米歇尔·\trAudin}}%
\newcommand{\trDamianFull}{\trn{Mihai Damian}{米哈伊·\trDamian}}%
\newcommand{\trZehnder}{\trn{Zehnder}{曾德尔}}%
\newcommand{\trMcDuff}{\trn{McDuff}{麦克达夫}}%
\newcommand{\trSalamon}{\trn{Salamon}{萨拉蒙}}%
\newcommand{\trHofer}{\trn{Hofer}{霍费尔}}%
\newcommand{\trEliashberg}{\trn{Eliashberg}{埃利亚什伯格}}%
\newcommand{\trWeinstein}{\trn{Weinstein}{温斯坦}}%
\newcommand{\trViterbo}{\trn{Viterbo}{维泰博}}%
\newcommand{\trSikorav}{\trn{Sikorav}{西科拉夫}}%
\newcommand{\trLaudenbach}{\trn{Laudenbach}{洛当巴克}}%
\newcommand{\trSchlenk}{\trn{Schlenk}{施伦克}}%
\newcommand{\trSchwarz}{\trn{Schwarz}{施瓦茨}}%
\newcommand{\trErne}{\trn{Ern\'e}{埃尔内}}%
\newcommand{\trKunneth}{\trn{K\"unneth}{屈内特}}%
\newcommand{\trBorsuk}{\trn{Borsuk}{博尔苏克}}%
\newcommand{\trUlam}{\trn{Ulam}{乌拉姆}}%
\newcommand{\trReeb}{\trn{Reeb}{雷布}}%
\newcommand{\trOno}{\trn{Ono}{小野薫}}%
\newcommand{\trFukaya}{\trn{Fukaya}{深谷贤治}}%
\newcommand{\trLiu}{\trn{Liu}{刘小博}}%
\newcommand{\trTian}{\trn{Tian}{田刚}}%
\newcommand{\trManifold}{流形}%
\newcommand{\trManifoldWithBoundary}{带边流形}%
\newcommand{\trCorneredBoundary}{带角边}%
\newcommand{\trSubmanifold}{子流形}%
\newcommand{\trHypersurface}{超曲面}%
\newcommand{\trTopology}{拓扑}%
\newcommand{\trSphere}{球面}%
\newcommand{\trBouquetOfSpheres}{球面束}%
\newcommand{\trMulticonnectionSpace}{多重连接空间}%
\newcommand{\trCircle}{圆周}%
\newcommand{\trTorus}{环面}%
\newcommand{\trDisk}{圆盘}%
\newcommand{\trSmooth}{光滑}%
\newcommand{\trChart}{坐标卡}%
\newcommand{\trAtlas}{图册}%
\newcommand{\trCriticalPoint}{临界点}%
\newcommand{\trCriticalValue}{临界值}%
\newcommand{\trRegularValue}{正则值}%
\newcommand{\trRegular}{正则}%
\newcommand{\trCriticalLevel}{临界水平}%
\newcommand{\trLevelSet}{水平集}%
\newcommand{\trSublevelSet}{次水平集}%
\newcommand{\trSaddlePoint}{鞍点}%
\newcommand{\trAntipodal}{对径}%
\newcommand{\trRectangularHyperbola}{等轴双曲线}%
\newcommand{\trBubble}{气泡}%
\newcommand{\trAsphericity}{非球面性}%
\newcommand{\trConic}{二次曲线}%
\newcommand{\trExtremum}{极值}%
\newcommand{\trHeightFunction}{高度函数}%
\newcommand{\trFocalPoint}{焦点}%
\newcommand{\trSingularPoint}{奇点}%
\newcommand{\trDegenerate}{退化}%
\newcommand{\trNondegenerate}{非退化}%
\newcommand{\trPositiveDefinite}{正定}%
\newcommand{\trNegativeDefinite}{负定}%
\newcommand{\trIsolated}{孤立}%
\newcommand{\trPlateauFunction}{平台函数}%
\newcommand{\trQuadric}{二次曲面}%
\newcommand{\trProjective}{射影}%
\newcommand{\trProjectiveSpace}{射影空间}%
\newcommand{\trProjectivePlane}{射影平面}%
\newcommand{\trHomogeneousCoordinates}{齐次坐标}%
\newcommand{\trHessian}{海塞}%
\newcommand{\trQuadraticForm}{二次型}%
\newcommand{\trSignature}{符号差}%
\newcommand{\trSymmetricBilinearForm}{对称双线性形式}%
\newcommand{\trBilinearForm}{双线性形式}%
\newcommand{\trWellDefined}{良定义}%
\newcommand{\trMorseFunction}{\trMorse{}函数}%
\newcommand{\trMorseLemma}{\trMorse{}引理}%
\newcommand{\trMorseChart}{\trMorse{}坐标卡}%
\newcommand{\trGradient}{梯度}%
\newcommand{\trVectorField}{向量场}%
\newcommand{\trVectorSpace}{向量空间}%
\newcommand{\trPseudogradient}{伪梯度}%
\newcommand{\trAdapted}{适配}%
\newcommand{\trIncoming}{内向}%
\newcommand{\trCollar}{领圈}%
\newcommand{\trExcision}{切除}%
\newcommand{\trStationary}{驻定}%
\newcommand{\trTrajectory}{轨道}%
\newcommand{\trFlowLine}{流线}%
\newcommand{\trIntegralCurve}{积分曲线}%
\newcommand{\trAutonomous}{自治}%
\newcommand{\trFlow}{流}%
\newcommand{\trStableManifold}{稳定流形}%
\newcommand{\trUnstableManifold}{不稳定流形}%
\newcommand{\trSmaleCondition}{\trSmale{}条件}%
\newcommand{\trIndex}{指标}%
\newcommand{\trCompact}{紧致}%
\newcommand{\trCompactification}{紧化}%
\newcommand{\trConnected}{连通}%
\newcommand{\trConnectedness}{连通性}%
\newcommand{\trConnectedComponent}{连通分支}%
\newcommand{\trSimplyConnected}{单连通}%
\newcommand{\trDense}{稠密}%
\newcommand{\trOpen}{开}%
\newcommand{\trGeneric}{通有}%
\newcommand{\trGenericity}{通有性}%
\newcommand{\trTransversality}{横截性}%
\newcommand{\trTransverse}{横截}%
\newcommand{\trSardTheorem}{\trSard{}定理}%
\newcommand{\trSubmersion}{浸没}%
\newcommand{\trImmersion}{浸入}%
\newcommand{\trEmbedding}{嵌入}%
\newcommand{\trParametrization}{参数化}%
\newcommand{\trDiffeomorphism}{微分同胚}%
\newcommand{\trDiffeomorphic}{微分同胚}%
\newcommand{\trHomeomorphism}{同胚}%
\newcommand{\trHomeomorphic}{同胚}%
\newcommand{\trTangentSpace}{切空间}%
\newcommand{\trTangentPlane}{切平面}%
\newcommand{\trTangentMap}{切映射}%
\newcommand{\trTangentBundle}{切丛}%
\newcommand{\trCotangent}{余切}%
\newcommand{\trVectorBundle}{向量丛}%
\newcommand{\trFiberBundle}{纤维丛}%
\newcommand{\trNormalBundle}{法丛}%
\newcommand{\trNormalVector}{法向量}%
\newcommand{\trTubularNeighborhood}{管状邻域}%
\newcommand{\trFiber}{纤维}%
\newcommand{\trSection}{截面}%
\newcommand{\trZeroSection}{零截面}%
\newcommand{\trTrivialization}{平凡化}%
\newcommand{\trMetric}{度量}%
\newcommand{\trRiemannian}{\trRiemann}%
\newcommand{\trEuclidean}{欧几里得}%
\newcommand{\trExponentialMap}{指数映射}%
\newcommand{\trExponentialDecay}{指数衰减}%
\newcommand{\trPartitionOfUnity}{单位分解}%
\newcommand{\trSupport}{支集}%
\newcommand{\trHandle}{手柄}%
\newcommand{\trCWComplex}{CW~复形}%
\newcommand{\trComplex}{复形}%
\newcommand{\trGenerator}{生成元}%
\newcommand{\trCardinality}{基数}%
\newcommand{\trDisjointUnion}{不交并}%
\newcommand{\trCommutativeRing}{交换环}%
\newcommand{\trAbelianGroup}{阿贝尔群}%
\newcommand{\trDegree}{次数}%
\newcommand{\trCellular}{胞腔}%
\newcommand{\trCellularDecomposition}{胞腔分解}%
\newcommand{\trSkeleton}{骨架}%
\newcommand{\trAttachingMap}{贴附映射}%
\newcommand{\trMorseComplex}{\trMorse{}复形}%
\newcommand{\trMorseHomology}{\trMorse{}同调}%
\newcommand{\trModuliSpace}{模空间}%
\newcommand{\trBrokenTrajectory}{破碎轨迹}%
\newcommand{\trGluing}{粘合}%
\newcommand{\trPregluing}{预粘合}%
\newcommand{\trQuotientTopology}{商拓扑}%
\newcommand{\trApproximateSolution}{近似解}%
\newcommand{\trExitPoint}{穿出点}%
\newcommand{\trEntryPoint}{穿入点}%
\newcommand{\trParallelTransport}{平行移动}%
\newcommand{\trSubbundle}{子丛}%
\newcommand{\trProper}{逆紧}%
\newcommand{\trLift}{提升}%
\newcommand{\trContractingMap}{压缩映射}%
\newcommand{\trSeparable}{可分}%
\newcommand{\trResolvent}{预解}%
\newcommand{\trHomology}{同调}%
\newcommand{\trCohomology}{上同调}%
\newcommand{\trHomologyGroup}{同调群}%
\newcommand{\trInvariant}{不变量}%
\newcommand{\trInvariance}{不变性}%
\newcommand{\trSingular}{奇异}%
\newcommand{\trNonsingular}{非奇异}%
\newcommand{\trSymplecticFrame}{辛标架}%
\newcommand{\trSymplecticDual}{辛对偶}%
\newcommand{\trSymplecticTransvection}{辛平延}%
\newcommand{\trCayley}{\trn{Cayley}{凯莱}}%
\newcommand{\trOctonion}{八元数}%
\newcommand{\trDerivation}{导子}%
\newcommand{\trParacompact}{仿紧}%
\newcommand{\trLie}{\trn{Lie}{李}}%
\newcommand{\trLieBracket}{\trLie{}括号}%
\newcommand{\trLieDerivative}{\trLie{}导数}%
\newcommand{\trInteriorProduct}{内乘}%
\newcommand{\trCovariantDerivative}{协变导数}%
\newcommand{\trConnection}{联络}%
\newcommand{\trGeodesic}{测地线}%
\newcommand{\trCartan}{\trn{Cartan}{嘉当}}%
\newcommand{\trEulerClass}{\trEuler{}类}%
\newcommand{\trAnnihilator}{零化子}%
\newcommand{\trConvolution}{卷积}%
\newcommand{\trFundamentalSolution}{基本解}%
\newcommand{\trWeakDerivative}{弱导数}%
\newcommand{\trHarmonic}{调和}%
\newcommand{\trLaplacian}{拉普拉斯算子}%
\newcommand{\trPoissonKernel}{\trPoisson{}核}%
\newcommand{\trIsotopy}{同痕}%
\newcommand{\trContact}{接触}%
\newcommand{\trSymplectization}{辛化}%
\newcommand{\trClosedCharacteristic}{闭特征}%
\newcommand{\trFrame}{标架}%
\newcommand{\trTrace}{迹}%
\newcommand{\trRelative}{相对}%
\newcommand{\trChainComplex}{链复形}%
\newcommand{\trBoundaryOperator}{边缘算子}%
\newcommand{\trDifferential}{微分}%
\newcommand{\trCellularHomology}{胞腔同调}%
\newcommand{\trBettiNumber}{贝蒂数}%
\newcommand{\trEulerCharacteristic}{欧拉示性数}%
\newcommand{\trPoincarePolynomial}{\trPoincare{}多项式}%
\newcommand{\trPoincareDuality}{\trPoincare{}对偶}%
\newcommand{\trKunnethFormula}{\trKunneth{}公式}%
\newcommand{\trExact}{正合}%
\newcommand{\trLongExactSequence}{长正合列}%
\newcommand{\trShortExactSequence}{短正合列}%
\newcommand{\trMorphism}{态射}%
\newcommand{\trHomomorphism}{同态}%
\newcommand{\trCanonical}{典范}%
\newcommand{\trConnectingHomomorphism}{连接同态}%
\newcommand{\trFunctor}{函子}%
\newcommand{\trFunctoriality}{函子性}%
\newcommand{\trCategory}{范畴}%
\newcommand{\trCobordism}{配边}%
\newcommand{\trOrientation}{定向}%
\newcommand{\trOrientable}{可定向}%
\newcommand{\trCoorientation}{余定向}%
\newcommand{\trComplement}{补空间}%
\newcommand{\trGraded}{分次}%
\newcommand{\trModule}{模}%
\newcommand{\trFundamentalGroup}{基本群}%
\newcommand{\trCovering}{覆盖}%
\newcommand{\trUniversalCover}{万有覆盖}%
\newcommand{\trHomotopy}{同伦}%
\newcommand{\trDeformationRetract}{形变收缩}%
\newcommand{\trInterpolation}{插值}%
\newcommand{\trDeformation}{形变}%
\newcommand{\trHomotopic}{同伦}%
\newcommand{\trMappingDegree}{映射度}%
\newcommand{\trIsomorphism}{同构}%
\newcommand{\trEndomorphism}{自同态}%
\newcommand{\trAutomorphism}{自同构}%
\newcommand{\trGeneralLinearGroup}{一般线性群}%
\newcommand{\trStabilizer}{稳定子群}%
\newcommand{\trHomogeneousSpace}{齐性空间}%
\newcommand{\trIsomorphic}{同构}%
\newcommand{\trKernel}{核}%
\newcommand{\trRank}{秩}%
\newcommand{\trInvertible}{可逆}%
\newcommand{\trCokernel}{余核}%
\newcommand{\trImage}{像}%
\newcommand{\trPreimage}{原像}%
\newcommand{\trInjective}{单射}%
\newcommand{\trSurjective}{满射}%
\newcommand{\trBijective}{双射}%
\newcommand{\trTensorProduct}{张量积}%
\newcommand{\trTorsion}{挠}%
\newcommand{\trDirectSum}{直和}%
\newcommand{\trLocalCoefficient}{局部系数}%
\newcommand{\trCoherentOrientation}{相容定向}%
\newcommand{\trSymplectic}{辛}%
\newcommand{\trSymplecticGroup}{辛群}%
\newcommand{\trSymplecticForm}{辛形式}%
\newcommand{\trSymplectomorphism}{辛同胚}%
\newcommand{\trSkewSymmetric}{斜对称}%
\newcommand{\trAlternating}{交错}%
\newcommand{\trAntisymmetric}{反对称}%
\newcommand{\trHermitian}{埃尔米特}%
\newcommand{\trKaehler}{凯勒}%
\newcommand{\trCalibrated}{标定}%
\newcommand{\trIsometry}{等距}%
\newcommand{\trEigenvalue}{特征值}%
\newcommand{\trCharacteristicPolynomial}{特征多项式}%
\newcommand{\trMultiplicity}{重数}%
\newcommand{\trEigenspace}{特征子空间}%
\newcommand{\trGeneralizedEigenspace}{广义特征子空间}%
\newcommand{\trDiagonalizable}{可对角化}%
\newcommand{\trComplexification}{复化}%
\newcommand{\trSymmetricMatrix}{对称矩阵}%
\newcommand{\trBlockDiagonalMatrix}{对角分块矩阵}%
\newcommand{\trMaximalCompactSubgroup}{极大紧子群}%
\newcommand{\trMonodromy}{单值映射}%
\newcommand{\trClosed}{闭}%
\newcommand{\trForm}{形式}%
\newcommand{\trVolume}{体积}%
\newcommand{\trAlmostComplex}{近复}%
\newcommand{\trComplexStructure}{复结构}%
\newcommand{\trCompatible}{相容}%
\newcommand{\trTamed}{驯服}%
\newcommand{\trHamiltonian}{\trHamilton}%
\newcommand{\trActionFunctional}{作用泛函}%
\newcommand{\trActionForm}{作用形式}%
\newcommand{\trIntegrationCover}{积分覆盖}%
\newcommand{\trInjectivityRadius}{单射半径}%
\newcommand{\trFrechet}{\trn{Fr\'echet}{弗雷歇}}%
\newcommand{\trGalois}{\trn{Galois}{伽罗瓦}}%
\newcommand{\trFloerHomology}{\trFloer{}同调}%
\newcommand{\trFloerComplex}{\trFloer{}复形}%
\newcommand{\trFloerEquation}{\trFloer{}方程}%
\newcommand{\trArnoldConjecture}{\trArnold{}猜想}%
\newcommand{\trMaslovIndex}{\trMaslov{}指标}%
\newcommand{\trConleyZehnderIndex}{\trConley--\trZehnder{}指标}%
\newcommand{\trMaslovCycle}{\trMaslov{}闭链}%
\newcommand{\trCycle}{闭链}%
\newcommand{\trLoop}{环路}%
\newcommand{\trLoopSpace}{环路空间}%
\newcommand{\trContractible}{可缩}%
\newcommand{\trPeriodic}{周期}%
\newcommand{\trFixedPoint}{不动点}%
\newcommand{\trCappingDisk}{帽盘}%
\newcommand{\trCappingSphere}{帽球面}%
\newcommand{\trCappingSurface}{帽曲面}%
\newcommand{\trChernClass}{\trChern{}类}%
\newcommand{\trFirstChernClass}{第一\trChern{}类}%
\newcommand{\trTautological}{重言}%
\newcommand{\trLagrangian}{\trLagrange{}式}%
\newcommand{\trLagrangianGrassmannian}{\trLagrange{}--\trGrassmann{}流形}%
\newcommand{\trIsotropic}{迷向}%
\newcommand{\trCoisotropic}{余迷向}%
\newcommand{\trLegendrian}{\trLegendre{}式}%
\newcommand{\trGeneratingFunction}{生成函数}%
\newcommand{\trUnitary}{酉}%
\newcommand{\trOrthogonal}{正交}%
\newcommand{\trOrthonormal}{标准正交}%
\newcommand{\trInnerProduct}{内积}%
\newcommand{\trPolarDecomposition}{极分解}%
\newcommand{\trSpectralFlow}{谱流}%
\newcommand{\trCrossing}{穿越点}%
\newcommand{\trPseudoholomorphicCurve}{伪全纯曲线}%
\newcommand{\trHolomorphic}{全纯}%
\newcommand{\trCauchyRiemann}{\trCauchy--\trRiemann}%
\newcommand{\trMonotone}{单调}%
\newcommand{\trEnergy}{能量}%
\newcommand{\trMomentum}{动量}%
\newcommand{\trEnergyIdentity}{能量恒等式}%
\newcommand{\trContinuationMorphism}{延续同态}%
\newcommand{\trSimilarityPrinciple}{相似性原理}%
\newcommand{\trContinuationPrinciple}{延续原理}%
\newcommand{\trMultiplePoint}{重点}%
\newcommand{\trInfiniteOrderZero}{无穷阶零点}%
\newcommand{\trSomewhereInjective}{某处单射}%
\newcommand{\trPerturbation}{扰动}%
\newcommand{\trLinearization}{线性化}%
\newcommand{\trBanachManifold}{\trBanach{}流形}%
\newcommand{\trFredholmOperator}{\trFredholm{}算子}%
\newcommand{\trBanachSpace}{\trBanach{}空间}%
\newcommand{\trHilbertSpace}{\trHilbert{}空间}%
\newcommand{\trSobolevSpace}{\trSobolev{}空间}%
\newcommand{\trWeakSolution}{弱解}%
\newcommand{\trWeakConvergence}{弱收敛}%
\newcommand{\trConvergence}{收敛}%
\newcommand{\trDistribution}{分布}%
\newcommand{\trElliptic}{椭圆}%
\newcommand{\trEllipticRegularity}{椭圆正则性}%
\newcommand{\trLinearOperator}{线性算子}%
\newcommand{\trBounded}{有界}%
\newcommand{\trSelfAdjoint}{自伴}%
\newcommand{\trAdjointOperator}{伴随算子}%
\newcommand{\trUnboundedOperator}{无界算子}%
\newcommand{\trHilbertianBasis}{希尔伯特基}%
\newcommand{\trSemiflow}{半流}%
\newcommand{\trSemigroup}{半群}%
\newcommand{\trSpectrum}{谱}%
\newcommand{\trUniformConvergence}{一致收敛}%
\newcommand{\trUniform}{一致}%
\newcommand{\trEquicontinuity}{等度连续}%
\newcommand{\trArzelaAscoliTheorem}{\trArzela--\trAscoli{}定理}%
\newcommand{\trImplicitFunctionTheorem}{隐函数定理}%
\newcommand{\trLocalInversionTheorem}{局部反演定理}%
\newcommand{\trRegularity}{正则性}%
\newcommand{\trApriori}{先验}%


\newcounter{secthm}[section]
\renewcommand{\thesecthm}{\thesection.\arabic{secthm}}
\newcounter{chex}[chapter]
\renewcommand{\thechex}{\arabic{chex}}
\newcommand{\vkthmsync}[1]{\setcounter{secthm}{\numexpr(#1)-1\relax}}
\newcommand{\vkexsync}[1]{\setcounter{chex}{\numexpr(#1)-1\relax}}

\newtheoremstyle{adplain}{\topsep}{\topsep}{\itshape}{}
  {\bfseries}{}{.5em}{#1~#2\ifstrempty{#3}{}{（{\kaishu #3}）}}
\newtheoremstyle{adup}{\topsep}{\topsep}{}{}
  {\bfseries}{}{.5em}{#1~#2\ifstrempty{#3}{}{（{\kaishu #3}）}}

\theoremstyle{adplain}
\newtheorem{theorem}    [secthm]{\trhTheorem}
\newtheorem{proposition}[secthm]{\trhProposition}
\newtheorem{lemma}      [secthm]{\trhLemma}
\newtheorem{corollary}  [secthm]{\trhCorollary}
\newtheorem{conjecture} [secthm]{\trhConjecture}
\newtheorem{assumption} [secthm]{\trhAssumption}
\theoremstyle{adup}
\newtheorem{definition} [secthm]{\trhDefinition}
\newtheorem{example}    [secthm]{\trhExample}
\newtheorem{notation}   [secthm]{\trhNotation}
\newtheorem{remark}     [secthm]{\trhRemark}
\newtheorem{warning}    [secthm]{\trhWarning}
\theoremstyle{adup}
\newtheorem{exercise}   [chex]  {\trhExercise}

\renewcommand{\proofname}{证明}
\makeatletter
\renewenvironment{proof}[1][\proofname]{\par
  \pushQED{\qed}\normalfont
  \topsep6\p@\@plus6\p@\relax
  \trivlist
  \item[\hskip\labelsep{\kaishu #1。}]\ignorespaces
}
  {\popQED\endtrivlist\@endpefalse}
\makeatother
\renewcommand{\qedsymbol}{$\vdash$}
\newcommand{\noqed}{\hfill$\vdash$\par}

\counterwithin{figure}{chapter}
\counterwithin{table}{chapter}

\newcommand{\epigraph}[2]{%
  \begin{center}\begin{minipage}{0.86\textwidth}
    {\kaishu #1}\par\vspace{2pt}
    \hfill{\kaishu\small #2}\par
  \end{minipage}\end{center}\vspace{6pt}
}

\newcommand{\RR}{\mathbf{R}}      %
\newcommand{\CC}{\mathbf{C}}
\newcommand{\ZZ}{\mathbf{Z}}
\newcommand{\NN}{\mathbf{N}}
\newcommand{\QQ}{\mathbf{Q}}
\newcommand{\cC}{\mathcal{C}}     %
\newcommand{\CP}{\mathbf{P}}      %
\newcommand{\Hom}{\operatorname{Hom}}
\newcommand{\id}{\mathrm{id}}
\newcommand{\im}{\operatorname{im}}
\newcommand{\Diff}{\operatorname{Diff}}
\newcommand{\scr}[1]{\mathscr{#1}}
\newcommand{\opp}{\mathrm{opp}}   %

\usepackage{xcolor}
\usepackage[hyperindex=false]{hyperref}
\usepackage{imakeidx}
\makeindex[title={译本术语索引},intoc,noautomatic]
\hypersetup{
  colorlinks=true,
  linkcolor=blue!55!black,
  citecolor=teal,
  urlcolor=violet,
  bookmarksnumbered=true,
  pdflang={zh-CN},
  pdftitle={莫尔斯理论与弗洛尔同调（中文译本）},
  pdfauthor={Michèle Audin, Mihai Damian（原著）},
}

\usepackage[capitalise,noabbrev]{cleveref}
\makeatletter
\let\mtfh@cref@linked\cref
\renewcommand{\cref}[1]{\mtfh@cref@linked*{#1}}
\makeatother
\crefname{theorem}    {\trhTheorem}    {\trhTheorem}
\crefname{proposition}{\trhProposition}{\trhProposition}
\crefname{lemma}      {\trhLemma}      {\trhLemma}
\crefname{corollary}  {\trhCorollary}  {\trhCorollary}
\crefname{conjecture} {\trhConjecture} {\trhConjecture}
\crefname{assumption} {\trhAssumption} {\trhAssumption}
\crefname{definition} {\trhDefinition} {\trhDefinition}
\crefname{example}    {\trhExample}    {\trhExample}
\crefname{notation}   {\trhNotation}   {\trhNotation}
\crefname{remark}     {\trhRemark}     {\trhRemark}
\crefname{warning}    {\trhWarning}    {\trhWarning}
\crefname{exercise}   {\trhExercise}   {\trhExercise}
\crefname{figure}{图}{图}
\crefname{table}{表}{表}
\crefformat{section}{§#2#1#3}
\crefformat{subsection}{§#2#1#3}
\crefname{secthm}{\trhTheorem}{\trhTheorem}
\crefname{chex}{\trhExercise}{\trhExercise}
\AtBeginEnvironment{theorem}    {\crefalias{secthm}{theorem}}
\AtBeginEnvironment{proposition}{\crefalias{secthm}{proposition}}
\AtBeginEnvironment{lemma}      {\crefalias{secthm}{lemma}}
\AtBeginEnvironment{corollary}  {\crefalias{secthm}{corollary}}
\AtBeginEnvironment{conjecture} {\crefalias{secthm}{conjecture}}
\AtBeginEnvironment{assumption} {\crefalias{secthm}{assumption}}
\AtBeginEnvironment{definition} {\crefalias{secthm}{definition}}
\AtBeginEnvironment{example}    {\crefalias{secthm}{example}}
\AtBeginEnvironment{notation}   {\crefalias{secthm}{notation}}
\AtBeginEnvironment{remark}     {\crefalias{secthm}{remark}}
\AtBeginEnvironment{warning}    {\crefalias{secthm}{warning}}
\AtBeginEnvironment{exercise}   {\crefalias{chex}{exercise}}

\newcommand{\BookGlossEntry}[2]{%
  \par\noindent\hangindent=1em\hangafter=1 #1\quad #2\par}
\newcommand{\BookSourceIndexEntry}[2]{%
  \par\begingroup\leftskip=#1em\relax\hangindent=1em\relax
  \noindent #2\par\endgroup}
\makeatletter
\newif\ifBookIndexing
\BookIndexingfalse
\newcommand{\BookBoldPage}[1]{\textbf{\hyperpage{#1}}}
\newcommand{\BookIndexDefinitionPage}[2]{%
  \expandafter\gdef\csname BookDefinition@#1@#2\endcsname{}}
\newcommand{\BookIndexPageFormatter}[2]{%
  \ifcsname BookDefinition@#1@#2\endcsname
    BookBoldPage%
  \else
    hyperpage%
  \fi}
\DeclareRobustCommand{\BookRecordIndex}[1]{%
  \ifBookIndexing
    \expandafter\index\expandafter{\csname BookIndex@#1\endcsname|%
      \noexpand\BookIndexPageFormatter{#1}{\thepage}}%
  \fi}
\DeclareRobustCommand{\BookIndexDefinition}[1]{%
  \ifBookIndexing
    \protected@write\@auxout{}{\string\BookIndexDefinitionPage{#1}{\thepage}}%
    \expandafter\index\expandafter{\csname BookIndex@#1\endcsname|BookBoldPage}%
  \fi}
\DeclareRobustCommand{\BookTerm}[2]{%
  \BookRecordIndex{#1}%
  \begingroup\let\BookRecordIndex\@gobble #2\endgroup}
\AddToHook{cmd/@outputpage/before}{\let\BookRecordIndex\@gobble}
\pdfstringdefDisableCommands{%
  \def\BookTerm#1#2{#2}%
  \let\BookRecordIndex\@gobble
  \let\BookIndexDefinition\@gobble}
\expandafter\def\csname BookIndex@trCalderonZygmund\endcsname{kaerdelongqigemengde@卡尔德龙--齐格蒙德（Calderón–Zygmund）}
\expandafter\def\csname BookIndex@trHilleYosida\endcsname{xileijitian@希勒--吉田（Hille–Yosida）}
\expandafter\def\csname BookIndex@trHahnBanach\endcsname{haenbanahe@哈恩--巴拿赫（Hahn–Banach）}
\expandafter\def\csname BookIndex@trManifold\endcsname{liuxing@流形（manifold）}
\expandafter\def\csname BookIndex@trManifoldWithBoundary\endcsname{daibianliuxing@带边流形（manifold with boundary）}
\expandafter\def\csname BookIndex@trCorneredBoundary\endcsname{daijiaobian@带角边（cornered boundary）}
\expandafter\def\csname BookIndex@trSubmanifold\endcsname{ziliuxing@子流形（submanifold）}
\expandafter\def\csname BookIndex@trHypersurface\endcsname{chaoqumian@超曲面（hypersurface）}
\expandafter\def\csname BookIndex@trTopology\endcsname{tuopu@拓扑（topology）}
\expandafter\def\csname BookIndex@trSphere\endcsname{qiumian@球面（sphere）}
\expandafter\def\csname BookIndex@trBouquetOfSpheres\endcsname{qiumianshu@球面束（bouquet of spheres）}
\expandafter\def\csname BookIndex@trMulticonnectionSpace\endcsname{duozhonglianjiekongjian@多重连接空间（multiconnection space）}
\expandafter\def\csname BookIndex@trCircle\endcsname{yuanzhou@圆周（circle）}
\expandafter\def\csname BookIndex@trTorus\endcsname{huanmian@环面（torus）}
\expandafter\def\csname BookIndex@trDisk\endcsname{yuanpan@圆盘（disk）}
\expandafter\def\csname BookIndex@trChart\endcsname{zuobiaoka@坐标卡（chart）}
\expandafter\def\csname BookIndex@trAtlas\endcsname{tuce@图册（atlas）}
\expandafter\def\csname BookIndex@trCriticalPoint\endcsname{linjiedian@临界点（critical point）}
\expandafter\def\csname BookIndex@trCriticalValue\endcsname{linjiezhi@临界值（critical value）}
\expandafter\def\csname BookIndex@trRegularValue\endcsname{zhengzezhi@正则值（regular value）}
\expandafter\def\csname BookIndex@trCriticalLevel\endcsname{linjieshuiping@临界水平（critical level）}
\expandafter\def\csname BookIndex@trLevelSet\endcsname{shuipingji@水平集（level set）}
\expandafter\def\csname BookIndex@trSublevelSet\endcsname{cishuipingji@次水平集（sublevel set）}
\expandafter\def\csname BookIndex@trSaddlePoint\endcsname{andian@鞍点（saddle point）}
\expandafter\def\csname BookIndex@trRectangularHyperbola\endcsname{dengzhoushuangquxian@等轴双曲线（rectangular hyperbola）}
\expandafter\def\csname BookIndex@trBubble\endcsname{qipao@气泡（bubble）}
\expandafter\def\csname BookIndex@trAsphericity\endcsname{feiqiumianxing@非球面性（asphericity）}
\expandafter\def\csname BookIndex@trConic\endcsname{erciquxian@二次曲线（conic）}
\expandafter\def\csname BookIndex@trExtremum\endcsname{jizhi@极值（extremum）}
\expandafter\def\csname BookIndex@trHeightFunction\endcsname{gaoduhanshu@高度函数（height function）}
\expandafter\def\csname BookIndex@trFocalPoint\endcsname{jiaodian@焦点（focal point）}
\expandafter\def\csname BookIndex@trSingularPoint\endcsname{qidian@奇点（singular point）}
\expandafter\def\csname BookIndex@trPlateauFunction\endcsname{pingtaihanshu@平台函数（plateau function）}
\expandafter\def\csname BookIndex@trQuadric\endcsname{erciqumian@二次曲面（quadric）}
\expandafter\def\csname BookIndex@trProjectiveSpace\endcsname{sheyingkongjian@射影空间（projective space）}
\expandafter\def\csname BookIndex@trProjectivePlane\endcsname{sheyingpingmian@射影平面（projective plane）}
\expandafter\def\csname BookIndex@trHomogeneousCoordinates\endcsname{qicizuobiao@齐次坐标（homogeneous coordinates）}
\expandafter\def\csname BookIndex@trHessian\endcsname{haisai@海塞（Hessian）}
\expandafter\def\csname BookIndex@trQuadraticForm\endcsname{ercixing@二次型（quadratic form）}
\expandafter\def\csname BookIndex@trSignature\endcsname{fuhaocha@符号差（signature）}
\expandafter\def\csname BookIndex@trSymmetricBilinearForm\endcsname{duichenshuangxianxingxingshi@对称双线性形式（symmetric bilinear form）}
\expandafter\def\csname BookIndex@trBilinearForm\endcsname{shuangxianxingxingshi@双线性形式（bilinear form）}
\expandafter\def\csname BookIndex@trMorseFunction\endcsname{moersihanshu@莫尔斯函数（Morse function）}
\expandafter\def\csname BookIndex@trMorseLemma\endcsname{moersiyinli@莫尔斯引理（Morse lemma）}
\expandafter\def\csname BookIndex@trMorseChart\endcsname{moersizuobiaoka@莫尔斯坐标卡（Morse chart）}
\expandafter\def\csname BookIndex@trGradient\endcsname{tidu@梯度（gradient）}
\expandafter\def\csname BookIndex@trVectorField\endcsname{xiangliangchang@向量场（vector field）}
\expandafter\def\csname BookIndex@trVectorSpace\endcsname{xiangliangkongjian@向量空间（vector space）}
\expandafter\def\csname BookIndex@trPseudogradient\endcsname{weitidu@伪梯度（pseudo-gradient）}
\expandafter\def\csname BookIndex@trCollar\endcsname{lingquan@领圈（collar）}
\expandafter\def\csname BookIndex@trExcision\endcsname{qiechu@切除（excision）}
\expandafter\def\csname BookIndex@trTrajectory\endcsname{guidao@轨道（trajectory / orbit）}
\expandafter\def\csname BookIndex@trFlowLine\endcsname{liuxian@流线（flow line）}
\expandafter\def\csname BookIndex@trIntegralCurve\endcsname{jifenquxian@积分曲线（integral curve）}
\expandafter\def\csname BookIndex@trFlow\endcsname{liu@流（flow）}
\expandafter\def\csname BookIndex@trStableManifold\endcsname{wendingliuxing@稳定流形（stable manifold）}
\expandafter\def\csname BookIndex@trUnstableManifold\endcsname{buwendingliuxing@不稳定流形（unstable manifold）}
\expandafter\def\csname BookIndex@trSmaleCondition\endcsname{simeiertiaojian@斯梅尔条件（Smale condition）}
\expandafter\def\csname BookIndex@trIndex\endcsname{zhibiao@指标（index）}
\expandafter\def\csname BookIndex@trCompact\endcsname{jinzhi@紧致（compact）}
\expandafter\def\csname BookIndex@trCompactification\endcsname{jinhua@紧化（compactification）}
\expandafter\def\csname BookIndex@trConnectedness\endcsname{liantongxing@连通性（connectedness）}
\expandafter\def\csname BookIndex@trConnectedComponent\endcsname{liantongfenzhi@连通分支（connected component）}
\expandafter\def\csname BookIndex@trGenericity\endcsname{tongyouxing@通有性（genericity）}
\expandafter\def\csname BookIndex@trTransversality\endcsname{hengjiexing@横截性（transversality）}
\expandafter\def\csname BookIndex@trSardTheorem\endcsname{sadedingli@萨德定理（Sard's theorem）}
\expandafter\def\csname BookIndex@trSubmersion\endcsname{jinmei@浸没（submersion）}
\expandafter\def\csname BookIndex@trImmersion\endcsname{jinru@浸入（immersion）}
\expandafter\def\csname BookIndex@trEmbedding\endcsname{qianru@嵌入（embedding）}
\expandafter\def\csname BookIndex@trParametrization\endcsname{canshuhua@参数化（parametrization）}
\expandafter\def\csname BookIndex@trDiffeomorphism\endcsname{weifentongpei@微分同胚（diffeomorphism）}
\expandafter\def\csname BookIndex@trHomeomorphism\endcsname{tongpei@同胚（homeomorphism）}
\expandafter\def\csname BookIndex@trTangentSpace\endcsname{qiekongjian@切空间（tangent space）}
\expandafter\def\csname BookIndex@trTangentPlane\endcsname{qiepingmian@切平面（tangent plane）}
\expandafter\def\csname BookIndex@trTangentMap\endcsname{qieyingshe@切映射（tangent map）}
\expandafter\def\csname BookIndex@trTangentBundle\endcsname{qiecong@切丛（tangent bundle）}
\expandafter\def\csname BookIndex@trCotangent\endcsname{yuqie@余切（cotangent）}
\expandafter\def\csname BookIndex@trVectorBundle\endcsname{xiangliangcong@向量丛（vector bundle）}
\expandafter\def\csname BookIndex@trFiberBundle\endcsname{xianweicong@纤维丛（fiber bundle）}
\expandafter\def\csname BookIndex@trNormalBundle\endcsname{facong@法丛（normal bundle）}
\expandafter\def\csname BookIndex@trNormalVector\endcsname{faxiangliang@法向量（normal vector）}
\expandafter\def\csname BookIndex@trTubularNeighborhood\endcsname{guanzhuanglinyu@管状邻域（tubular neighborhood）}
\expandafter\def\csname BookIndex@trFiber\endcsname{xianwei@纤维（fiber）}
\expandafter\def\csname BookIndex@trSection\endcsname{jiemian@截面（section）}
\expandafter\def\csname BookIndex@trZeroSection\endcsname{lingjiemian@零截面（zero section）}
\expandafter\def\csname BookIndex@trTrivialization\endcsname{pingfanhua@平凡化（trivialization）}
\expandafter\def\csname BookIndex@trMetric\endcsname{duliang@度量（metric）}
\expandafter\def\csname BookIndex@trExponentialMap\endcsname{zhishuyingshe@指数映射（exponential map）}
\expandafter\def\csname BookIndex@trExponentialDecay\endcsname{zhishushuaijian@指数衰减（exponential decay）}
\expandafter\def\csname BookIndex@trPartitionOfUnity\endcsname{danweifenjie@单位分解（partition of unity）}
\expandafter\def\csname BookIndex@trSupport\endcsname{zhiji@支集（support）}
\expandafter\def\csname BookIndex@trHandle\endcsname{shoubing@手柄（handle）}
\expandafter\def\csname BookIndex@trCWComplex\endcsname{cwfuxing@CW~复形（CW-complex）}
\expandafter\def\csname BookIndex@trComplex\endcsname{fuxing@复形（complex）}
\expandafter\def\csname BookIndex@trGenerator\endcsname{shengchengyuan@生成元（generator）}
\expandafter\def\csname BookIndex@trCardinality\endcsname{jishu@基数（cardinality）}
\expandafter\def\csname BookIndex@trDisjointUnion\endcsname{bujiaobing@不交并（disjoint union）}
\expandafter\def\csname BookIndex@trCommutativeRing\endcsname{jiaohuanhuan@交换环（commutative ring）}
\expandafter\def\csname BookIndex@trAbelianGroup\endcsname{abeierqun@阿贝尔群（abelian group）}
\expandafter\def\csname BookIndex@trDegree\endcsname{cishu@次数（degree）}
\expandafter\def\csname BookIndex@trCellularDecomposition\endcsname{baoqiangfenjie@胞腔分解（cellular decomposition）}
\expandafter\def\csname BookIndex@trSkeleton\endcsname{gujia@骨架（skeleton）}
\expandafter\def\csname BookIndex@trAttachingMap\endcsname{tiefuyingshe@贴附映射（attaching map）}
\expandafter\def\csname BookIndex@trMorseComplex\endcsname{moersifuxing@莫尔斯复形（Morse complex）}
\expandafter\def\csname BookIndex@trMorseHomology\endcsname{moersitongdiao@莫尔斯同调（Morse homology）}
\expandafter\def\csname BookIndex@trModuliSpace\endcsname{mokongjian@模空间（moduli space）}
\expandafter\def\csname BookIndex@trBrokenTrajectory\endcsname{posuiguiji@破碎轨迹（broken trajectory）}
\expandafter\def\csname BookIndex@trGluing\endcsname{zhanhe@粘合（gluing）}
\expandafter\def\csname BookIndex@trPregluing\endcsname{yuzhanhe@预粘合（pre-gluing）}
\expandafter\def\csname BookIndex@trQuotientTopology\endcsname{shangtuopu@商拓扑（quotient topology）}
\expandafter\def\csname BookIndex@trApproximateSolution\endcsname{jinsijie@近似解（approximate solution）}
\expandafter\def\csname BookIndex@trExitPoint\endcsname{chuanchudian@穿出点（exit point）}
\expandafter\def\csname BookIndex@trEntryPoint\endcsname{chuanrudian@穿入点（entry point）}
\expandafter\def\csname BookIndex@trParallelTransport\endcsname{pingxingyidong@平行移动（parallel transport）}
\expandafter\def\csname BookIndex@trSubbundle\endcsname{zicong@子丛（subbundle）}
\expandafter\def\csname BookIndex@trLift\endcsname{tisheng@提升（lift）}
\expandafter\def\csname BookIndex@trContractingMap\endcsname{yasuoyingshe@压缩映射（contracting map）}
\expandafter\def\csname BookIndex@trResolvent\endcsname{yujie@预解（resolvent）}
\expandafter\def\csname BookIndex@trHomology\endcsname{tongdiao@同调（homology）}
\expandafter\def\csname BookIndex@trCohomology\endcsname{shangtongdiao@上同调（cohomology）}
\expandafter\def\csname BookIndex@trHomologyGroup\endcsname{tongdiaoqun@同调群（homology group）}
\expandafter\def\csname BookIndex@trInvariant\endcsname{bubianliang@不变量（invariant）}
\expandafter\def\csname BookIndex@trInvariance\endcsname{bubianxing@不变性（invariance）}
\expandafter\def\csname BookIndex@trSymplecticFrame\endcsname{xinbiaojia@辛标架（symplectic frame）}
\expandafter\def\csname BookIndex@trSymplecticDual\endcsname{xinduiou@辛对偶（symplectic dual）}
\expandafter\def\csname BookIndex@trSymplecticTransvection\endcsname{xinpingyan@辛平延（symplectic transvection）}
\expandafter\def\csname BookIndex@trOctonion\endcsname{bayuanshu@八元数（octaves／octonions）}
\expandafter\def\csname BookIndex@trDerivation\endcsname{daozi@导子（derivation）}
\expandafter\def\csname BookIndex@trLieBracket\endcsname{likuohao@李括号（Lie bracket）}
\expandafter\def\csname BookIndex@trLieDerivative\endcsname{lidaoshu@李导数（Lie derivative）}
\expandafter\def\csname BookIndex@trInteriorProduct\endcsname{neicheng@内乘（interior product）}
\expandafter\def\csname BookIndex@trCovariantDerivative\endcsname{xiebiandaoshu@协变导数（covariant derivative）}
\expandafter\def\csname BookIndex@trConnection\endcsname{lianluo@联络（connection）}
\expandafter\def\csname BookIndex@trGeodesic\endcsname{cedexian@测地线（geodesic）}
\expandafter\def\csname BookIndex@trEulerClass\endcsname{oulalei@欧拉类（Euler class）}
\expandafter\def\csname BookIndex@trAnnihilator\endcsname{linghuazi@零化子（annihilator）}
\expandafter\def\csname BookIndex@trConvolution\endcsname{juanji@卷积（convolution）}
\expandafter\def\csname BookIndex@trFundamentalSolution\endcsname{jibenjie@基本解（fundamental solution）}
\expandafter\def\csname BookIndex@trWeakDerivative\endcsname{ruodaoshu@弱导数（weak derivative）}
\expandafter\def\csname BookIndex@trLaplacian\endcsname{lapulasisuanzi@拉普拉斯算子（Laplacian）}
\expandafter\def\csname BookIndex@trPoissonKernel\endcsname{posonghe@泊松核（Poisson kernel）}
\expandafter\def\csname BookIndex@trIsotopy\endcsname{tonghen@同痕（isotopy）}
\expandafter\def\csname BookIndex@trContact\endcsname{jiechu@接触（contact）}
\expandafter\def\csname BookIndex@trSymplectization\endcsname{xinhua@辛化（symplectization）}
\expandafter\def\csname BookIndex@trClosedCharacteristic\endcsname{bitezheng@闭特征（closed characteristic）}
\expandafter\def\csname BookIndex@trFrame\endcsname{biaojia@标架（frame）}
\expandafter\def\csname BookIndex@trTrace\endcsname{ji@迹（trace）}
\expandafter\def\csname BookIndex@trChainComplex\endcsname{lianfuxing@链复形（chain complex）}
\expandafter\def\csname BookIndex@trBoundaryOperator\endcsname{bianyuansuanzi@边缘算子（boundary operator）}
\expandafter\def\csname BookIndex@trDifferential\endcsname{weifen@微分（differential）}
\expandafter\def\csname BookIndex@trCellularHomology\endcsname{baoqiangtongdiao@胞腔同调（cellular homology）}
\expandafter\def\csname BookIndex@trBettiNumber\endcsname{beidishu@贝蒂数（Betti number）}
\expandafter\def\csname BookIndex@trEulerCharacteristic\endcsname{oulashixingshu@欧拉示性数（Euler characteristic）}
\expandafter\def\csname BookIndex@trPoincarePolynomial\endcsname{pangjialaiduoxiangshi@庞加莱多项式（Poincaré polynomial）}
\expandafter\def\csname BookIndex@trPoincareDuality\endcsname{pangjialaiduiou@庞加莱对偶（Poincaré duality）}
\expandafter\def\csname BookIndex@trKunnethFormula\endcsname{kunnethgongshi@屈内特公式（Künneth formula）}
\expandafter\def\csname BookIndex@trLongExactSequence\endcsname{changzhenghelie@长正合列（long exact sequence）}
\expandafter\def\csname BookIndex@trShortExactSequence\endcsname{duanzhenghelie@短正合列（short exact sequence）}
\expandafter\def\csname BookIndex@trMorphism\endcsname{taishe@态射（morphism）}
\expandafter\def\csname BookIndex@trHomomorphism\endcsname{tongtai@同态（homomorphism）}
\expandafter\def\csname BookIndex@trConnectingHomomorphism\endcsname{lianjietongtai@连接同态（connecting homomorphism）}
\expandafter\def\csname BookIndex@trFunctor\endcsname{hanzi@函子（functor）}
\expandafter\def\csname BookIndex@trFunctoriality\endcsname{hanzixing@函子性（functoriality）}
\expandafter\def\csname BookIndex@trCategory\endcsname{fanchou@范畴（category）}
\expandafter\def\csname BookIndex@trCobordism\endcsname{peibian@配边（cobordism）}
\expandafter\def\csname BookIndex@trOrientation\endcsname{dingxiang@定向（orientation）}
\expandafter\def\csname BookIndex@trCoorientation\endcsname{yudingxiang@余定向（co-orientation / co-oriented）}
\expandafter\def\csname BookIndex@trComplement\endcsname{bukongjian@补空间（complement）}
\expandafter\def\csname BookIndex@trModule\endcsname{mo@模（module）}
\expandafter\def\csname BookIndex@trFundamentalGroup\endcsname{jibenqun@基本群（fundamental group）}
\expandafter\def\csname BookIndex@trCovering\endcsname{fugai@覆盖（covering）}
\expandafter\def\csname BookIndex@trUniversalCover\endcsname{wanyoufugai@万有覆盖（universal cover）}
\expandafter\def\csname BookIndex@trHomotopy\endcsname{tonglun@同伦（homotopy）}
\expandafter\def\csname BookIndex@trDeformationRetract\endcsname{xingbianshousuo@形变收缩（deformation retract）}
\expandafter\def\csname BookIndex@trInterpolation\endcsname{chazhi@插值（interpolation）}
\expandafter\def\csname BookIndex@trDeformation\endcsname{xingbian@形变（deformation）}
\expandafter\def\csname BookIndex@trMappingDegree\endcsname{yingshedu@映射度（mapping degree）}
\expandafter\def\csname BookIndex@trIsomorphism\endcsname{tonggou@同构（isomorphism）}
\expandafter\def\csname BookIndex@trEndomorphism\endcsname{zitongtai@自同态（endomorphism）}
\expandafter\def\csname BookIndex@trAutomorphism\endcsname{zitonggou@自同构（automorphism）}
\expandafter\def\csname BookIndex@trGeneralLinearGroup\endcsname{yibanxianxingqun@一般线性群（general linear group）}
\expandafter\def\csname BookIndex@trStabilizer\endcsname{wendingziqun@稳定子群（stabilizer）}
\expandafter\def\csname BookIndex@trHomogeneousSpace\endcsname{qixingkongjian@齐性空间（homogeneous space）}
\expandafter\def\csname BookIndex@trKernel\endcsname{he@核（kernel）}
\expandafter\def\csname BookIndex@trRank\endcsname{zhi@秩（rank）}
\expandafter\def\csname BookIndex@trCokernel\endcsname{yuhe@余核（cokernel）}
\expandafter\def\csname BookIndex@trImage\endcsname{xiang@像（image）}
\expandafter\def\csname BookIndex@trPreimage\endcsname{yuanxiang@原像（preimage）}
\expandafter\def\csname BookIndex@trTensorProduct\endcsname{zhangliangji@张量积（tensor product）}
\expandafter\def\csname BookIndex@trTorsion\endcsname{nao@挠（torsion）}
\expandafter\def\csname BookIndex@trDirectSum\endcsname{zhihe@直和（direct sum）}
\expandafter\def\csname BookIndex@trLocalCoefficient\endcsname{jubuxishu@局部系数（local coefficients）}
\expandafter\def\csname BookIndex@trCoherentOrientation\endcsname{xiangrongdingxiang@相容定向（coherent orientation）}
\expandafter\def\csname BookIndex@trSymplecticGroup\endcsname{xinqun@辛群（symplectic group）}
\expandafter\def\csname BookIndex@trSymplecticForm\endcsname{xinxingshi@辛形式（symplectic form）}
\expandafter\def\csname BookIndex@trSymplectomorphism\endcsname{xintongpei@辛同胚（symplectomorphism）}
\expandafter\def\csname BookIndex@trIsometry\endcsname{dengju@等距（isometry）}
\expandafter\def\csname BookIndex@trEigenvalue\endcsname{tezhengzhi@特征值（eigenvalue）}
\expandafter\def\csname BookIndex@trCharacteristicPolynomial\endcsname{tezhengduoxiangshi@特征多项式（characteristic polynomial）}
\expandafter\def\csname BookIndex@trMultiplicity\endcsname{chongshu@重数（multiplicity）}
\expandafter\def\csname BookIndex@trEigenspace\endcsname{tezhengzikongjian@特征子空间（eigenspace）}
\expandafter\def\csname BookIndex@trGeneralizedEigenspace\endcsname{guangyitezhengzikongjian@广义特征子空间（generalized eigenspace）}
\expandafter\def\csname BookIndex@trComplexification\endcsname{fuhua@复化（complexification）}
\expandafter\def\csname BookIndex@trSymmetricMatrix\endcsname{duichenjuzhen@对称矩阵（symmetric matrix）}
\expandafter\def\csname BookIndex@trBlockDiagonalMatrix\endcsname{duijiaofenkuaijuzhen@对角分块矩阵（block-diagonal matrix）}
\expandafter\def\csname BookIndex@trMaximalCompactSubgroup\endcsname{jidajinziqun@极大紧子群（maximal compact subgroup）}
\expandafter\def\csname BookIndex@trMonodromy\endcsname{danzhiyingshe@单值映射（monodromy）}
\expandafter\def\csname BookIndex@trForm\endcsname{xingshi@形式（form）}
\expandafter\def\csname BookIndex@trVolume\endcsname{tiji@体积（volume）}
\expandafter\def\csname BookIndex@trComplexStructure\endcsname{fujiegou@复结构（complex structure）}
\expandafter\def\csname BookIndex@trActionFunctional\endcsname{zuoyongfanhan@作用泛函（action functional）}
\expandafter\def\csname BookIndex@trActionForm\endcsname{zuoyongxingshi@作用形式（action form）}
\expandafter\def\csname BookIndex@trIntegrationCover\endcsname{jifenfugai@积分覆盖（integration cover）}
\expandafter\def\csname BookIndex@trInjectivityRadius\endcsname{danshebanjing@单射半径（injectivity radius）}
\expandafter\def\csname BookIndex@trFloerHomology\endcsname{fuluoertongdiao@弗洛尔同调（Floer homology）}
\expandafter\def\csname BookIndex@trFloerComplex\endcsname{fuluoerfuxing@弗洛尔复形（Floer complex）}
\expandafter\def\csname BookIndex@trFloerEquation\endcsname{fuluoerfangcheng@弗洛尔方程（Floer equation）}
\expandafter\def\csname BookIndex@trArnoldConjecture\endcsname{anuodecaixiang@阿诺德猜想（Arnold conjecture）}
\expandafter\def\csname BookIndex@trMaslovIndex\endcsname{masiluofuzhibiao@马斯洛夫指标（Maslov index）}
\expandafter\def\csname BookIndex@trConleyZehnderIndex\endcsname{kanglizehnderzhibiao@康利--曾德尔指标（Conley–Zehnder index）}
\expandafter\def\csname BookIndex@trMaslovCycle\endcsname{masiluofubilian@马斯洛夫闭链（Maslov cycle）}
\expandafter\def\csname BookIndex@trCycle\endcsname{bilian@闭链（cycle）}
\expandafter\def\csname BookIndex@trLoop\endcsname{huanlu@环路（loop）}
\expandafter\def\csname BookIndex@trLoopSpace\endcsname{huanlukongjian@环路空间（loop space）}
\expandafter\def\csname BookIndex@trFixedPoint\endcsname{budongdian@不动点（fixed point）}
\expandafter\def\csname BookIndex@trCappingDisk\endcsname{maopan@帽盘（capping disk）}
\expandafter\def\csname BookIndex@trCappingSphere\endcsname{maoqiumian@帽球面（capping sphere）}
\expandafter\def\csname BookIndex@trCappingSurface\endcsname{maoqumian@帽曲面（capping surface）}
\expandafter\def\csname BookIndex@trChernClass\endcsname{chenlei@陈类（Chern class）}
\expandafter\def\csname BookIndex@trFirstChernClass\endcsname{diyichenlei@第一陈类（first Chern class）}
\expandafter\def\csname BookIndex@trLagrangianGrassmannian\endcsname{lagelangrigelasimanliuxing@拉格朗日--格拉斯曼流形（Lagrangian Grassmannian）}
\expandafter\def\csname BookIndex@trGeneratingFunction\endcsname{shengchenghanshu@生成函数（generating function）}
\expandafter\def\csname BookIndex@trInnerProduct\endcsname{neiji@内积（inner product）}
\expandafter\def\csname BookIndex@trPolarDecomposition\endcsname{jifenjie@极分解（polar decomposition）}
\expandafter\def\csname BookIndex@trSpectralFlow\endcsname{puliu@谱流（spectral flow）}
\expandafter\def\csname BookIndex@trCrossing\endcsname{chuanyuedian@穿越点（crossing）}
\expandafter\def\csname BookIndex@trPseudoholomorphicCurve\endcsname{weiquanchunquxian@伪全纯曲线（pseudoholomorphic curve）}
\expandafter\def\csname BookIndex@trCauchyRiemann\endcsname{kexiliman@柯西--黎曼（Cauchy–Riemann）}
\expandafter\def\csname BookIndex@trEnergy\endcsname{nengliang@能量（energy）}
\expandafter\def\csname BookIndex@trMomentum\endcsname{dongliang@动量（momentum）}
\expandafter\def\csname BookIndex@trEnergyIdentity\endcsname{nenglianghengdengshi@能量恒等式（energy identity）}
\expandafter\def\csname BookIndex@trContinuationMorphism\endcsname{yanxutongtai@延续同态（continuation morphism）}
\expandafter\def\csname BookIndex@trSimilarityPrinciple\endcsname{xiangsixingyuanli@相似性原理（similarity principle）}
\expandafter\def\csname BookIndex@trContinuationPrinciple\endcsname{yanxuyuanli@延续原理（continuation principle）}
\expandafter\def\csname BookIndex@trMultiplePoint\endcsname{chongdian@重点（multiple point）}
\expandafter\def\csname BookIndex@trInfiniteOrderZero\endcsname{wuqiongjielingdian@无穷阶零点（infinite-order zero）}
\expandafter\def\csname BookIndex@trPerturbation\endcsname{raodong@扰动（perturbation）}
\expandafter\def\csname BookIndex@trLinearization\endcsname{xianxinghua@线性化（linearization）}
\expandafter\def\csname BookIndex@trBanachManifold\endcsname{banaheliuxing@巴拿赫流形（Banach manifold）}
\expandafter\def\csname BookIndex@trFredholmOperator\endcsname{fuleidehuomusuanzi@弗雷德霍姆算子（Fredholm operator）}
\expandafter\def\csname BookIndex@trBanachSpace\endcsname{banahekongjian@巴拿赫空间（Banach space）}
\expandafter\def\csname BookIndex@trHilbertSpace\endcsname{xierbotekongjian@希尔伯特空间（Hilbert space）}
\expandafter\def\csname BookIndex@trSobolevSpace\endcsname{suoboliefukongjian@索伯列夫空间（Sobolev space）}
\expandafter\def\csname BookIndex@trWeakSolution\endcsname{ruojie@弱解（weak solution）}
\expandafter\def\csname BookIndex@trWeakConvergence\endcsname{ruoshoulian@弱收敛（weak convergence）}
\expandafter\def\csname BookIndex@trConvergence\endcsname{shoulian@收敛（convergence）}
\expandafter\def\csname BookIndex@trDistribution\endcsname{fenbu@分布（distribution）}
\expandafter\def\csname BookIndex@trEllipticRegularity\endcsname{tuoyuanzhengzexing@椭圆正则性（elliptic regularity）}
\expandafter\def\csname BookIndex@trLinearOperator\endcsname{xianxingsuanzi@线性算子（linear operator）}
\expandafter\def\csname BookIndex@trAdjointOperator\endcsname{bansuisuanzi@伴随算子（adjoint operator）}
\expandafter\def\csname BookIndex@trUnboundedOperator\endcsname{wujiesuanzi@无界算子（unbounded operator）}
\expandafter\def\csname BookIndex@trHilbertianBasis\endcsname{xierboteji@希尔伯特基（Hilbertian basis）}
\expandafter\def\csname BookIndex@trSemiflow\endcsname{banliu@半流（semiflow）}
\expandafter\def\csname BookIndex@trSemigroup\endcsname{banqun@半群（semigroup）}
\expandafter\def\csname BookIndex@trSpectrum\endcsname{pu@谱（spectrum）}
\expandafter\def\csname BookIndex@trUniformConvergence\endcsname{yizhishoulian@一致收敛（uniform convergence）}
\expandafter\def\csname BookIndex@trEquicontinuity\endcsname{dengdulianxu@等度连续（equicontinuity）}
\expandafter\def\csname BookIndex@trArzelaAscoliTheorem\endcsname{aerzelaasikelidingli@阿尔泽拉--阿斯科利定理（Arzelà–Ascoli theorem）}
\expandafter\def\csname BookIndex@trImplicitFunctionTheorem\endcsname{yinhanshudingli@隐函数定理（implicit function theorem）}
\expandafter\def\csname BookIndex@trLocalInversionTheorem\endcsname{jubufanyandingli@局部反演定理（local inversion theorem）}
\expandafter\def\csname BookIndex@trRegularity\endcsname{zhengzexing@正则性（regularity）}
\expandafter\def\csname BookIndex@T001\endcsname{xinliuxing@辛流形（symplectic manifold）}
\expandafter\def\csname BookIndex@T002\endcsname{xinxiangliangkongjian@辛向量空间（symplectic vector space）}
\expandafter\def\csname BookIndex@T003\endcsname{xinji@辛基（symplectic basis）}
\expandafter\def\csname BookIndex@T004\endcsname{dabudingli@达布定理（Darboux theorem）}
\expandafter\def\csname BookIndex@T005\endcsname{lagelangriziliuxing@拉格朗日子流形（Lagrangian submanifold）}
\expandafter\def\csname BookIndex@T006\endcsname{hamidunweifentongpeiqun@哈密顿微分同胚群（Hamiltonian diffeomorphism group）}
\expandafter\def\csname BookIndex@T007\endcsname{jiechuliuxing@接触流形（contact manifold）}
\expandafter\def\csname BookIndex@T008\endcsname{moersibudengshi@莫尔斯不等式（Morse inequalities）}
\expandafter\def\csname BookIndex@T009\endcsname{leibudingli@雷布定理（Reeb theorem）}
\expandafter\def\csname BookIndex@T010\endcsname{simeierdingli@斯梅尔定理（Smale theorem）}
\expandafter\def\csname BookIndex@T011\endcsname{kupukasimeierdingli@库普卡–斯梅尔定理（Kupka–Smale theorem）}
\expandafter\def\csname BookIndex@T012\endcsname{sadesimeierdingli@萨德–斯梅尔定理（Sard–Smale theorem）}
\expandafter\def\csname BookIndex@T013\endcsname{huiteniqianrudingli@惠特尼嵌入定理（Whitney embedding theorem）}
\expandafter\def\csname BookIndex@T014\endcsname{bulaoweierbudongdiandingli@布劳威尔不动点定理（Brouwer fixed point theorem）}
\expandafter\def\csname BookIndex@T015\endcsname{xileijitiandingli@希勒–吉田定理（Hille–Yosida theorem）}
\expandafter\def\csname BookIndex@T016\endcsname{leilixijinxingdingli@雷利希紧性定理（Rellich compactness theorem）}
\expandafter\def\csname BookIndex@T017\endcsname{leilixikongdelashaofudingli@雷利希–孔德拉绍夫定理（Rellich–Kondrachov theorem）}
\expandafter\def\csname BookIndex@T018\endcsname{fuleidehuomuzhibiaowendingxing@弗雷德霍姆指标稳定性（Fredholm index stability）}
\expandafter\let\csname BookOriginal@trCalderonZygmund\endcsname\trCalderonZygmund
\renewcommand{\trCalderonZygmund}{\BookTerm{trCalderonZygmund}{\csname BookOriginal@trCalderonZygmund\endcsname}}
\expandafter\let\csname BookOriginal@trHilleYosida\endcsname\trHilleYosida
\renewcommand{\trHilleYosida}{\BookTerm{trHilleYosida}{\csname BookOriginal@trHilleYosida\endcsname}}
\expandafter\let\csname BookOriginal@trHahnBanach\endcsname\trHahnBanach
\renewcommand{\trHahnBanach}{\BookTerm{trHahnBanach}{\csname BookOriginal@trHahnBanach\endcsname}}
\expandafter\let\csname BookOriginal@trManifold\endcsname\trManifold
\renewcommand{\trManifold}{\BookTerm{trManifold}{\csname BookOriginal@trManifold\endcsname}}
\expandafter\let\csname BookOriginal@trManifoldWithBoundary\endcsname\trManifoldWithBoundary
\renewcommand{\trManifoldWithBoundary}{\BookTerm{trManifoldWithBoundary}{\csname BookOriginal@trManifoldWithBoundary\endcsname}}
\expandafter\let\csname BookOriginal@trCorneredBoundary\endcsname\trCorneredBoundary
\renewcommand{\trCorneredBoundary}{\BookTerm{trCorneredBoundary}{\csname BookOriginal@trCorneredBoundary\endcsname}}
\expandafter\let\csname BookOriginal@trSubmanifold\endcsname\trSubmanifold
\renewcommand{\trSubmanifold}{\BookTerm{trSubmanifold}{\csname BookOriginal@trSubmanifold\endcsname}}
\expandafter\let\csname BookOriginal@trHypersurface\endcsname\trHypersurface
\renewcommand{\trHypersurface}{\BookTerm{trHypersurface}{\csname BookOriginal@trHypersurface\endcsname}}
\expandafter\let\csname BookOriginal@trTopology\endcsname\trTopology
\renewcommand{\trTopology}{\BookTerm{trTopology}{\csname BookOriginal@trTopology\endcsname}}
\expandafter\let\csname BookOriginal@trSphere\endcsname\trSphere
\renewcommand{\trSphere}{\BookTerm{trSphere}{\csname BookOriginal@trSphere\endcsname}}
\expandafter\let\csname BookOriginal@trBouquetOfSpheres\endcsname\trBouquetOfSpheres
\renewcommand{\trBouquetOfSpheres}{\BookTerm{trBouquetOfSpheres}{\csname BookOriginal@trBouquetOfSpheres\endcsname}}
\expandafter\let\csname BookOriginal@trMulticonnectionSpace\endcsname\trMulticonnectionSpace
\renewcommand{\trMulticonnectionSpace}{\BookTerm{trMulticonnectionSpace}{\csname BookOriginal@trMulticonnectionSpace\endcsname}}
\expandafter\let\csname BookOriginal@trCircle\endcsname\trCircle
\renewcommand{\trCircle}{\BookTerm{trCircle}{\csname BookOriginal@trCircle\endcsname}}
\expandafter\let\csname BookOriginal@trTorus\endcsname\trTorus
\renewcommand{\trTorus}{\BookTerm{trTorus}{\csname BookOriginal@trTorus\endcsname}}
\expandafter\let\csname BookOriginal@trDisk\endcsname\trDisk
\renewcommand{\trDisk}{\BookTerm{trDisk}{\csname BookOriginal@trDisk\endcsname}}
\expandafter\let\csname BookOriginal@trChart\endcsname\trChart
\renewcommand{\trChart}{\BookTerm{trChart}{\csname BookOriginal@trChart\endcsname}}
\expandafter\let\csname BookOriginal@trAtlas\endcsname\trAtlas
\renewcommand{\trAtlas}{\BookTerm{trAtlas}{\csname BookOriginal@trAtlas\endcsname}}
\expandafter\let\csname BookOriginal@trCriticalPoint\endcsname\trCriticalPoint
\renewcommand{\trCriticalPoint}{\BookTerm{trCriticalPoint}{\csname BookOriginal@trCriticalPoint\endcsname}}
\expandafter\let\csname BookOriginal@trCriticalValue\endcsname\trCriticalValue
\renewcommand{\trCriticalValue}{\BookTerm{trCriticalValue}{\csname BookOriginal@trCriticalValue\endcsname}}
\expandafter\let\csname BookOriginal@trRegularValue\endcsname\trRegularValue
\renewcommand{\trRegularValue}{\BookTerm{trRegularValue}{\csname BookOriginal@trRegularValue\endcsname}}
\expandafter\let\csname BookOriginal@trCriticalLevel\endcsname\trCriticalLevel
\renewcommand{\trCriticalLevel}{\BookTerm{trCriticalLevel}{\csname BookOriginal@trCriticalLevel\endcsname}}
\expandafter\let\csname BookOriginal@trLevelSet\endcsname\trLevelSet
\renewcommand{\trLevelSet}{\BookTerm{trLevelSet}{\csname BookOriginal@trLevelSet\endcsname}}
\expandafter\let\csname BookOriginal@trSublevelSet\endcsname\trSublevelSet
\renewcommand{\trSublevelSet}{\BookTerm{trSublevelSet}{\csname BookOriginal@trSublevelSet\endcsname}}
\expandafter\let\csname BookOriginal@trSaddlePoint\endcsname\trSaddlePoint
\renewcommand{\trSaddlePoint}{\BookTerm{trSaddlePoint}{\csname BookOriginal@trSaddlePoint\endcsname}}
\expandafter\let\csname BookOriginal@trRectangularHyperbola\endcsname\trRectangularHyperbola
\renewcommand{\trRectangularHyperbola}{\BookTerm{trRectangularHyperbola}{\csname BookOriginal@trRectangularHyperbola\endcsname}}
\expandafter\let\csname BookOriginal@trBubble\endcsname\trBubble
\renewcommand{\trBubble}{\BookTerm{trBubble}{\csname BookOriginal@trBubble\endcsname}}
\expandafter\let\csname BookOriginal@trAsphericity\endcsname\trAsphericity
\renewcommand{\trAsphericity}{\BookTerm{trAsphericity}{\csname BookOriginal@trAsphericity\endcsname}}
\expandafter\let\csname BookOriginal@trConic\endcsname\trConic
\renewcommand{\trConic}{\BookTerm{trConic}{\csname BookOriginal@trConic\endcsname}}
\expandafter\let\csname BookOriginal@trExtremum\endcsname\trExtremum
\renewcommand{\trExtremum}{\BookTerm{trExtremum}{\csname BookOriginal@trExtremum\endcsname}}
\expandafter\let\csname BookOriginal@trHeightFunction\endcsname\trHeightFunction
\renewcommand{\trHeightFunction}{\BookTerm{trHeightFunction}{\csname BookOriginal@trHeightFunction\endcsname}}
\expandafter\let\csname BookOriginal@trFocalPoint\endcsname\trFocalPoint
\renewcommand{\trFocalPoint}{\BookTerm{trFocalPoint}{\csname BookOriginal@trFocalPoint\endcsname}}
\expandafter\let\csname BookOriginal@trSingularPoint\endcsname\trSingularPoint
\renewcommand{\trSingularPoint}{\BookTerm{trSingularPoint}{\csname BookOriginal@trSingularPoint\endcsname}}
\expandafter\let\csname BookOriginal@trPlateauFunction\endcsname\trPlateauFunction
\renewcommand{\trPlateauFunction}{\BookTerm{trPlateauFunction}{\csname BookOriginal@trPlateauFunction\endcsname}}
\expandafter\let\csname BookOriginal@trQuadric\endcsname\trQuadric
\renewcommand{\trQuadric}{\BookTerm{trQuadric}{\csname BookOriginal@trQuadric\endcsname}}
\expandafter\let\csname BookOriginal@trProjectiveSpace\endcsname\trProjectiveSpace
\renewcommand{\trProjectiveSpace}{\BookTerm{trProjectiveSpace}{\csname BookOriginal@trProjectiveSpace\endcsname}}
\expandafter\let\csname BookOriginal@trProjectivePlane\endcsname\trProjectivePlane
\renewcommand{\trProjectivePlane}{\BookTerm{trProjectivePlane}{\csname BookOriginal@trProjectivePlane\endcsname}}
\expandafter\let\csname BookOriginal@trHomogeneousCoordinates\endcsname\trHomogeneousCoordinates
\renewcommand{\trHomogeneousCoordinates}{\BookTerm{trHomogeneousCoordinates}{\csname BookOriginal@trHomogeneousCoordinates\endcsname}}
\expandafter\let\csname BookOriginal@trHessian\endcsname\trHessian
\renewcommand{\trHessian}{\BookTerm{trHessian}{\csname BookOriginal@trHessian\endcsname}}
\expandafter\let\csname BookOriginal@trQuadraticForm\endcsname\trQuadraticForm
\renewcommand{\trQuadraticForm}{\BookTerm{trQuadraticForm}{\csname BookOriginal@trQuadraticForm\endcsname}}
\expandafter\let\csname BookOriginal@trSignature\endcsname\trSignature
\renewcommand{\trSignature}{\BookTerm{trSignature}{\csname BookOriginal@trSignature\endcsname}}
\expandafter\let\csname BookOriginal@trSymmetricBilinearForm\endcsname\trSymmetricBilinearForm
\renewcommand{\trSymmetricBilinearForm}{\BookTerm{trSymmetricBilinearForm}{\csname BookOriginal@trSymmetricBilinearForm\endcsname}}
\expandafter\let\csname BookOriginal@trBilinearForm\endcsname\trBilinearForm
\renewcommand{\trBilinearForm}{\BookTerm{trBilinearForm}{\csname BookOriginal@trBilinearForm\endcsname}}
\expandafter\let\csname BookOriginal@trMorseFunction\endcsname\trMorseFunction
\renewcommand{\trMorseFunction}{\BookTerm{trMorseFunction}{\csname BookOriginal@trMorseFunction\endcsname}}
\expandafter\let\csname BookOriginal@trMorseLemma\endcsname\trMorseLemma
\renewcommand{\trMorseLemma}{\BookTerm{trMorseLemma}{\csname BookOriginal@trMorseLemma\endcsname}}
\expandafter\let\csname BookOriginal@trMorseChart\endcsname\trMorseChart
\renewcommand{\trMorseChart}{\BookTerm{trMorseChart}{\csname BookOriginal@trMorseChart\endcsname}}
\expandafter\let\csname BookOriginal@trGradient\endcsname\trGradient
\renewcommand{\trGradient}{\BookTerm{trGradient}{\csname BookOriginal@trGradient\endcsname}}
\expandafter\let\csname BookOriginal@trVectorField\endcsname\trVectorField
\renewcommand{\trVectorField}{\BookTerm{trVectorField}{\csname BookOriginal@trVectorField\endcsname}}
\expandafter\let\csname BookOriginal@trVectorSpace\endcsname\trVectorSpace
\renewcommand{\trVectorSpace}{\BookTerm{trVectorSpace}{\csname BookOriginal@trVectorSpace\endcsname}}
\expandafter\let\csname BookOriginal@trPseudogradient\endcsname\trPseudogradient
\renewcommand{\trPseudogradient}{\BookTerm{trPseudogradient}{\csname BookOriginal@trPseudogradient\endcsname}}
\expandafter\let\csname BookOriginal@trCollar\endcsname\trCollar
\renewcommand{\trCollar}{\BookTerm{trCollar}{\csname BookOriginal@trCollar\endcsname}}
\expandafter\let\csname BookOriginal@trExcision\endcsname\trExcision
\renewcommand{\trExcision}{\BookTerm{trExcision}{\csname BookOriginal@trExcision\endcsname}}
\expandafter\let\csname BookOriginal@trTrajectory\endcsname\trTrajectory
\renewcommand{\trTrajectory}{\BookTerm{trTrajectory}{\csname BookOriginal@trTrajectory\endcsname}}
\expandafter\let\csname BookOriginal@trFlowLine\endcsname\trFlowLine
\renewcommand{\trFlowLine}{\BookTerm{trFlowLine}{\csname BookOriginal@trFlowLine\endcsname}}
\expandafter\let\csname BookOriginal@trIntegralCurve\endcsname\trIntegralCurve
\renewcommand{\trIntegralCurve}{\BookTerm{trIntegralCurve}{\csname BookOriginal@trIntegralCurve\endcsname}}
\expandafter\let\csname BookOriginal@trFlow\endcsname\trFlow
\renewcommand{\trFlow}{\BookTerm{trFlow}{\csname BookOriginal@trFlow\endcsname}}
\expandafter\let\csname BookOriginal@trStableManifold\endcsname\trStableManifold
\renewcommand{\trStableManifold}{\BookTerm{trStableManifold}{\csname BookOriginal@trStableManifold\endcsname}}
\expandafter\let\csname BookOriginal@trUnstableManifold\endcsname\trUnstableManifold
\renewcommand{\trUnstableManifold}{\BookTerm{trUnstableManifold}{\csname BookOriginal@trUnstableManifold\endcsname}}
\expandafter\let\csname BookOriginal@trSmaleCondition\endcsname\trSmaleCondition
\renewcommand{\trSmaleCondition}{\BookTerm{trSmaleCondition}{\csname BookOriginal@trSmaleCondition\endcsname}}
\expandafter\let\csname BookOriginal@trIndex\endcsname\trIndex
\renewcommand{\trIndex}{\BookTerm{trIndex}{\csname BookOriginal@trIndex\endcsname}}
\expandafter\let\csname BookOriginal@trCompact\endcsname\trCompact
\renewcommand{\trCompact}{\BookTerm{trCompact}{\csname BookOriginal@trCompact\endcsname}}
\expandafter\let\csname BookOriginal@trCompactification\endcsname\trCompactification
\renewcommand{\trCompactification}{\BookTerm{trCompactification}{\csname BookOriginal@trCompactification\endcsname}}
\expandafter\let\csname BookOriginal@trConnectedness\endcsname\trConnectedness
\renewcommand{\trConnectedness}{\BookTerm{trConnectedness}{\csname BookOriginal@trConnectedness\endcsname}}
\expandafter\let\csname BookOriginal@trConnectedComponent\endcsname\trConnectedComponent
\renewcommand{\trConnectedComponent}{\BookTerm{trConnectedComponent}{\csname BookOriginal@trConnectedComponent\endcsname}}
\expandafter\let\csname BookOriginal@trGenericity\endcsname\trGenericity
\renewcommand{\trGenericity}{\BookTerm{trGenericity}{\csname BookOriginal@trGenericity\endcsname}}
\expandafter\let\csname BookOriginal@trTransversality\endcsname\trTransversality
\renewcommand{\trTransversality}{\BookTerm{trTransversality}{\csname BookOriginal@trTransversality\endcsname}}
\expandafter\let\csname BookOriginal@trSardTheorem\endcsname\trSardTheorem
\renewcommand{\trSardTheorem}{\BookTerm{trSardTheorem}{\csname BookOriginal@trSardTheorem\endcsname}}
\expandafter\let\csname BookOriginal@trSubmersion\endcsname\trSubmersion
\renewcommand{\trSubmersion}{\BookTerm{trSubmersion}{\csname BookOriginal@trSubmersion\endcsname}}
\expandafter\let\csname BookOriginal@trImmersion\endcsname\trImmersion
\renewcommand{\trImmersion}{\BookTerm{trImmersion}{\csname BookOriginal@trImmersion\endcsname}}
\expandafter\let\csname BookOriginal@trEmbedding\endcsname\trEmbedding
\renewcommand{\trEmbedding}{\BookTerm{trEmbedding}{\csname BookOriginal@trEmbedding\endcsname}}
\expandafter\let\csname BookOriginal@trParametrization\endcsname\trParametrization
\renewcommand{\trParametrization}{\BookTerm{trParametrization}{\csname BookOriginal@trParametrization\endcsname}}
\expandafter\let\csname BookOriginal@trDiffeomorphism\endcsname\trDiffeomorphism
\renewcommand{\trDiffeomorphism}{\BookTerm{trDiffeomorphism}{\csname BookOriginal@trDiffeomorphism\endcsname}}
\expandafter\let\csname BookOriginal@trHomeomorphism\endcsname\trHomeomorphism
\renewcommand{\trHomeomorphism}{\BookTerm{trHomeomorphism}{\csname BookOriginal@trHomeomorphism\endcsname}}
\expandafter\let\csname BookOriginal@trTangentSpace\endcsname\trTangentSpace
\renewcommand{\trTangentSpace}{\BookTerm{trTangentSpace}{\csname BookOriginal@trTangentSpace\endcsname}}
\expandafter\let\csname BookOriginal@trTangentPlane\endcsname\trTangentPlane
\renewcommand{\trTangentPlane}{\BookTerm{trTangentPlane}{\csname BookOriginal@trTangentPlane\endcsname}}
\expandafter\let\csname BookOriginal@trTangentMap\endcsname\trTangentMap
\renewcommand{\trTangentMap}{\BookTerm{trTangentMap}{\csname BookOriginal@trTangentMap\endcsname}}
\expandafter\let\csname BookOriginal@trTangentBundle\endcsname\trTangentBundle
\renewcommand{\trTangentBundle}{\BookTerm{trTangentBundle}{\csname BookOriginal@trTangentBundle\endcsname}}
\expandafter\let\csname BookOriginal@trCotangent\endcsname\trCotangent
\renewcommand{\trCotangent}{\BookTerm{trCotangent}{\csname BookOriginal@trCotangent\endcsname}}
\expandafter\let\csname BookOriginal@trVectorBundle\endcsname\trVectorBundle
\renewcommand{\trVectorBundle}{\BookTerm{trVectorBundle}{\csname BookOriginal@trVectorBundle\endcsname}}
\expandafter\let\csname BookOriginal@trFiberBundle\endcsname\trFiberBundle
\renewcommand{\trFiberBundle}{\BookTerm{trFiberBundle}{\csname BookOriginal@trFiberBundle\endcsname}}
\expandafter\let\csname BookOriginal@trNormalBundle\endcsname\trNormalBundle
\renewcommand{\trNormalBundle}{\BookTerm{trNormalBundle}{\csname BookOriginal@trNormalBundle\endcsname}}
\expandafter\let\csname BookOriginal@trNormalVector\endcsname\trNormalVector
\renewcommand{\trNormalVector}{\BookTerm{trNormalVector}{\csname BookOriginal@trNormalVector\endcsname}}
\expandafter\let\csname BookOriginal@trTubularNeighborhood\endcsname\trTubularNeighborhood
\renewcommand{\trTubularNeighborhood}{\BookTerm{trTubularNeighborhood}{\csname BookOriginal@trTubularNeighborhood\endcsname}}
\expandafter\let\csname BookOriginal@trFiber\endcsname\trFiber
\renewcommand{\trFiber}{\BookTerm{trFiber}{\csname BookOriginal@trFiber\endcsname}}
\expandafter\let\csname BookOriginal@trSection\endcsname\trSection
\renewcommand{\trSection}{\BookTerm{trSection}{\csname BookOriginal@trSection\endcsname}}
\expandafter\let\csname BookOriginal@trZeroSection\endcsname\trZeroSection
\renewcommand{\trZeroSection}{\BookTerm{trZeroSection}{\csname BookOriginal@trZeroSection\endcsname}}
\expandafter\let\csname BookOriginal@trTrivialization\endcsname\trTrivialization
\renewcommand{\trTrivialization}{\BookTerm{trTrivialization}{\csname BookOriginal@trTrivialization\endcsname}}
\expandafter\let\csname BookOriginal@trMetric\endcsname\trMetric
\renewcommand{\trMetric}{\BookTerm{trMetric}{\csname BookOriginal@trMetric\endcsname}}
\expandafter\let\csname BookOriginal@trExponentialMap\endcsname\trExponentialMap
\renewcommand{\trExponentialMap}{\BookTerm{trExponentialMap}{\csname BookOriginal@trExponentialMap\endcsname}}
\expandafter\let\csname BookOriginal@trExponentialDecay\endcsname\trExponentialDecay
\renewcommand{\trExponentialDecay}{\BookTerm{trExponentialDecay}{\csname BookOriginal@trExponentialDecay\endcsname}}
\expandafter\let\csname BookOriginal@trPartitionOfUnity\endcsname\trPartitionOfUnity
\renewcommand{\trPartitionOfUnity}{\BookTerm{trPartitionOfUnity}{\csname BookOriginal@trPartitionOfUnity\endcsname}}
\expandafter\let\csname BookOriginal@trSupport\endcsname\trSupport
\renewcommand{\trSupport}{\BookTerm{trSupport}{\csname BookOriginal@trSupport\endcsname}}
\expandafter\let\csname BookOriginal@trHandle\endcsname\trHandle
\renewcommand{\trHandle}{\BookTerm{trHandle}{\csname BookOriginal@trHandle\endcsname}}
\expandafter\let\csname BookOriginal@trCWComplex\endcsname\trCWComplex
\renewcommand{\trCWComplex}{\BookTerm{trCWComplex}{\csname BookOriginal@trCWComplex\endcsname}}
\expandafter\let\csname BookOriginal@trComplex\endcsname\trComplex
\renewcommand{\trComplex}{\BookTerm{trComplex}{\csname BookOriginal@trComplex\endcsname}}
\expandafter\let\csname BookOriginal@trGenerator\endcsname\trGenerator
\renewcommand{\trGenerator}{\BookTerm{trGenerator}{\csname BookOriginal@trGenerator\endcsname}}
\expandafter\let\csname BookOriginal@trCardinality\endcsname\trCardinality
\renewcommand{\trCardinality}{\BookTerm{trCardinality}{\csname BookOriginal@trCardinality\endcsname}}
\expandafter\let\csname BookOriginal@trDisjointUnion\endcsname\trDisjointUnion
\renewcommand{\trDisjointUnion}{\BookTerm{trDisjointUnion}{\csname BookOriginal@trDisjointUnion\endcsname}}
\expandafter\let\csname BookOriginal@trCommutativeRing\endcsname\trCommutativeRing
\renewcommand{\trCommutativeRing}{\BookTerm{trCommutativeRing}{\csname BookOriginal@trCommutativeRing\endcsname}}
\expandafter\let\csname BookOriginal@trAbelianGroup\endcsname\trAbelianGroup
\renewcommand{\trAbelianGroup}{\BookTerm{trAbelianGroup}{\csname BookOriginal@trAbelianGroup\endcsname}}
\expandafter\let\csname BookOriginal@trDegree\endcsname\trDegree
\renewcommand{\trDegree}{\BookTerm{trDegree}{\csname BookOriginal@trDegree\endcsname}}
\expandafter\let\csname BookOriginal@trCellularDecomposition\endcsname\trCellularDecomposition
\renewcommand{\trCellularDecomposition}{\BookTerm{trCellularDecomposition}{\csname BookOriginal@trCellularDecomposition\endcsname}}
\expandafter\let\csname BookOriginal@trSkeleton\endcsname\trSkeleton
\renewcommand{\trSkeleton}{\BookTerm{trSkeleton}{\csname BookOriginal@trSkeleton\endcsname}}
\expandafter\let\csname BookOriginal@trAttachingMap\endcsname\trAttachingMap
\renewcommand{\trAttachingMap}{\BookTerm{trAttachingMap}{\csname BookOriginal@trAttachingMap\endcsname}}
\expandafter\let\csname BookOriginal@trMorseComplex\endcsname\trMorseComplex
\renewcommand{\trMorseComplex}{\BookTerm{trMorseComplex}{\csname BookOriginal@trMorseComplex\endcsname}}
\expandafter\let\csname BookOriginal@trMorseHomology\endcsname\trMorseHomology
\renewcommand{\trMorseHomology}{\BookTerm{trMorseHomology}{\csname BookOriginal@trMorseHomology\endcsname}}
\expandafter\let\csname BookOriginal@trModuliSpace\endcsname\trModuliSpace
\renewcommand{\trModuliSpace}{\BookTerm{trModuliSpace}{\csname BookOriginal@trModuliSpace\endcsname}}
\expandafter\let\csname BookOriginal@trBrokenTrajectory\endcsname\trBrokenTrajectory
\renewcommand{\trBrokenTrajectory}{\BookTerm{trBrokenTrajectory}{\csname BookOriginal@trBrokenTrajectory\endcsname}}
\expandafter\let\csname BookOriginal@trGluing\endcsname\trGluing
\renewcommand{\trGluing}{\BookTerm{trGluing}{\csname BookOriginal@trGluing\endcsname}}
\expandafter\let\csname BookOriginal@trPregluing\endcsname\trPregluing
\renewcommand{\trPregluing}{\BookTerm{trPregluing}{\csname BookOriginal@trPregluing\endcsname}}
\expandafter\let\csname BookOriginal@trQuotientTopology\endcsname\trQuotientTopology
\renewcommand{\trQuotientTopology}{\BookTerm{trQuotientTopology}{\csname BookOriginal@trQuotientTopology\endcsname}}
\expandafter\let\csname BookOriginal@trApproximateSolution\endcsname\trApproximateSolution
\renewcommand{\trApproximateSolution}{\BookTerm{trApproximateSolution}{\csname BookOriginal@trApproximateSolution\endcsname}}
\expandafter\let\csname BookOriginal@trExitPoint\endcsname\trExitPoint
\renewcommand{\trExitPoint}{\BookTerm{trExitPoint}{\csname BookOriginal@trExitPoint\endcsname}}
\expandafter\let\csname BookOriginal@trEntryPoint\endcsname\trEntryPoint
\renewcommand{\trEntryPoint}{\BookTerm{trEntryPoint}{\csname BookOriginal@trEntryPoint\endcsname}}
\expandafter\let\csname BookOriginal@trParallelTransport\endcsname\trParallelTransport
\renewcommand{\trParallelTransport}{\BookTerm{trParallelTransport}{\csname BookOriginal@trParallelTransport\endcsname}}
\expandafter\let\csname BookOriginal@trSubbundle\endcsname\trSubbundle
\renewcommand{\trSubbundle}{\BookTerm{trSubbundle}{\csname BookOriginal@trSubbundle\endcsname}}
\expandafter\let\csname BookOriginal@trLift\endcsname\trLift
\renewcommand{\trLift}{\BookTerm{trLift}{\csname BookOriginal@trLift\endcsname}}
\expandafter\let\csname BookOriginal@trContractingMap\endcsname\trContractingMap
\renewcommand{\trContractingMap}{\BookTerm{trContractingMap}{\csname BookOriginal@trContractingMap\endcsname}}
\expandafter\let\csname BookOriginal@trResolvent\endcsname\trResolvent
\renewcommand{\trResolvent}{\BookTerm{trResolvent}{\csname BookOriginal@trResolvent\endcsname}}
\expandafter\let\csname BookOriginal@trHomology\endcsname\trHomology
\renewcommand{\trHomology}{\BookTerm{trHomology}{\csname BookOriginal@trHomology\endcsname}}
\expandafter\let\csname BookOriginal@trCohomology\endcsname\trCohomology
\renewcommand{\trCohomology}{\BookTerm{trCohomology}{\csname BookOriginal@trCohomology\endcsname}}
\expandafter\let\csname BookOriginal@trHomologyGroup\endcsname\trHomologyGroup
\renewcommand{\trHomologyGroup}{\BookTerm{trHomologyGroup}{\csname BookOriginal@trHomologyGroup\endcsname}}
\expandafter\let\csname BookOriginal@trInvariant\endcsname\trInvariant
\renewcommand{\trInvariant}{\BookTerm{trInvariant}{\csname BookOriginal@trInvariant\endcsname}}
\expandafter\let\csname BookOriginal@trInvariance\endcsname\trInvariance
\renewcommand{\trInvariance}{\BookTerm{trInvariance}{\csname BookOriginal@trInvariance\endcsname}}
\expandafter\let\csname BookOriginal@trSymplecticFrame\endcsname\trSymplecticFrame
\renewcommand{\trSymplecticFrame}{\BookTerm{trSymplecticFrame}{\csname BookOriginal@trSymplecticFrame\endcsname}}
\expandafter\let\csname BookOriginal@trSymplecticDual\endcsname\trSymplecticDual
\renewcommand{\trSymplecticDual}{\BookTerm{trSymplecticDual}{\csname BookOriginal@trSymplecticDual\endcsname}}
\expandafter\let\csname BookOriginal@trSymplecticTransvection\endcsname\trSymplecticTransvection
\renewcommand{\trSymplecticTransvection}{\BookTerm{trSymplecticTransvection}{\csname BookOriginal@trSymplecticTransvection\endcsname}}
\expandafter\let\csname BookOriginal@trOctonion\endcsname\trOctonion
\renewcommand{\trOctonion}{\BookTerm{trOctonion}{\csname BookOriginal@trOctonion\endcsname}}
\expandafter\let\csname BookOriginal@trDerivation\endcsname\trDerivation
\renewcommand{\trDerivation}{\BookTerm{trDerivation}{\csname BookOriginal@trDerivation\endcsname}}
\expandafter\let\csname BookOriginal@trLieBracket\endcsname\trLieBracket
\renewcommand{\trLieBracket}{\BookTerm{trLieBracket}{\csname BookOriginal@trLieBracket\endcsname}}
\expandafter\let\csname BookOriginal@trLieDerivative\endcsname\trLieDerivative
\renewcommand{\trLieDerivative}{\BookTerm{trLieDerivative}{\csname BookOriginal@trLieDerivative\endcsname}}
\expandafter\let\csname BookOriginal@trInteriorProduct\endcsname\trInteriorProduct
\renewcommand{\trInteriorProduct}{\BookTerm{trInteriorProduct}{\csname BookOriginal@trInteriorProduct\endcsname}}
\expandafter\let\csname BookOriginal@trCovariantDerivative\endcsname\trCovariantDerivative
\renewcommand{\trCovariantDerivative}{\BookTerm{trCovariantDerivative}{\csname BookOriginal@trCovariantDerivative\endcsname}}
\expandafter\let\csname BookOriginal@trConnection\endcsname\trConnection
\renewcommand{\trConnection}{\BookTerm{trConnection}{\csname BookOriginal@trConnection\endcsname}}
\expandafter\let\csname BookOriginal@trGeodesic\endcsname\trGeodesic
\renewcommand{\trGeodesic}{\BookTerm{trGeodesic}{\csname BookOriginal@trGeodesic\endcsname}}
\expandafter\let\csname BookOriginal@trEulerClass\endcsname\trEulerClass
\renewcommand{\trEulerClass}{\BookTerm{trEulerClass}{\csname BookOriginal@trEulerClass\endcsname}}
\expandafter\let\csname BookOriginal@trAnnihilator\endcsname\trAnnihilator
\renewcommand{\trAnnihilator}{\BookTerm{trAnnihilator}{\csname BookOriginal@trAnnihilator\endcsname}}
\expandafter\let\csname BookOriginal@trConvolution\endcsname\trConvolution
\renewcommand{\trConvolution}{\BookTerm{trConvolution}{\csname BookOriginal@trConvolution\endcsname}}
\expandafter\let\csname BookOriginal@trFundamentalSolution\endcsname\trFundamentalSolution
\renewcommand{\trFundamentalSolution}{\BookTerm{trFundamentalSolution}{\csname BookOriginal@trFundamentalSolution\endcsname}}
\expandafter\let\csname BookOriginal@trWeakDerivative\endcsname\trWeakDerivative
\renewcommand{\trWeakDerivative}{\BookTerm{trWeakDerivative}{\csname BookOriginal@trWeakDerivative\endcsname}}
\expandafter\let\csname BookOriginal@trLaplacian\endcsname\trLaplacian
\renewcommand{\trLaplacian}{\BookTerm{trLaplacian}{\csname BookOriginal@trLaplacian\endcsname}}
\expandafter\let\csname BookOriginal@trPoissonKernel\endcsname\trPoissonKernel
\renewcommand{\trPoissonKernel}{\BookTerm{trPoissonKernel}{\csname BookOriginal@trPoissonKernel\endcsname}}
\expandafter\let\csname BookOriginal@trIsotopy\endcsname\trIsotopy
\renewcommand{\trIsotopy}{\BookTerm{trIsotopy}{\csname BookOriginal@trIsotopy\endcsname}}
\expandafter\let\csname BookOriginal@trContact\endcsname\trContact
\renewcommand{\trContact}{\BookTerm{trContact}{\csname BookOriginal@trContact\endcsname}}
\expandafter\let\csname BookOriginal@trSymplectization\endcsname\trSymplectization
\renewcommand{\trSymplectization}{\BookTerm{trSymplectization}{\csname BookOriginal@trSymplectization\endcsname}}
\expandafter\let\csname BookOriginal@trClosedCharacteristic\endcsname\trClosedCharacteristic
\renewcommand{\trClosedCharacteristic}{\BookTerm{trClosedCharacteristic}{\csname BookOriginal@trClosedCharacteristic\endcsname}}
\expandafter\let\csname BookOriginal@trFrame\endcsname\trFrame
\renewcommand{\trFrame}{\BookTerm{trFrame}{\csname BookOriginal@trFrame\endcsname}}
\expandafter\let\csname BookOriginal@trTrace\endcsname\trTrace
\renewcommand{\trTrace}{\BookTerm{trTrace}{\csname BookOriginal@trTrace\endcsname}}
\expandafter\let\csname BookOriginal@trChainComplex\endcsname\trChainComplex
\renewcommand{\trChainComplex}{\BookTerm{trChainComplex}{\csname BookOriginal@trChainComplex\endcsname}}
\expandafter\let\csname BookOriginal@trBoundaryOperator\endcsname\trBoundaryOperator
\renewcommand{\trBoundaryOperator}{\BookTerm{trBoundaryOperator}{\csname BookOriginal@trBoundaryOperator\endcsname}}
\expandafter\let\csname BookOriginal@trDifferential\endcsname\trDifferential
\renewcommand{\trDifferential}{\BookTerm{trDifferential}{\csname BookOriginal@trDifferential\endcsname}}
\expandafter\let\csname BookOriginal@trCellularHomology\endcsname\trCellularHomology
\renewcommand{\trCellularHomology}{\BookTerm{trCellularHomology}{\csname BookOriginal@trCellularHomology\endcsname}}
\expandafter\let\csname BookOriginal@trBettiNumber\endcsname\trBettiNumber
\renewcommand{\trBettiNumber}{\BookTerm{trBettiNumber}{\csname BookOriginal@trBettiNumber\endcsname}}
\expandafter\let\csname BookOriginal@trEulerCharacteristic\endcsname\trEulerCharacteristic
\renewcommand{\trEulerCharacteristic}{\BookTerm{trEulerCharacteristic}{\csname BookOriginal@trEulerCharacteristic\endcsname}}
\expandafter\let\csname BookOriginal@trPoincarePolynomial\endcsname\trPoincarePolynomial
\renewcommand{\trPoincarePolynomial}{\BookTerm{trPoincarePolynomial}{\csname BookOriginal@trPoincarePolynomial\endcsname}}
\expandafter\let\csname BookOriginal@trPoincareDuality\endcsname\trPoincareDuality
\renewcommand{\trPoincareDuality}{\BookTerm{trPoincareDuality}{\csname BookOriginal@trPoincareDuality\endcsname}}
\expandafter\let\csname BookOriginal@trKunnethFormula\endcsname\trKunnethFormula
\renewcommand{\trKunnethFormula}{\BookTerm{trKunnethFormula}{\csname BookOriginal@trKunnethFormula\endcsname}}
\expandafter\let\csname BookOriginal@trLongExactSequence\endcsname\trLongExactSequence
\renewcommand{\trLongExactSequence}{\BookTerm{trLongExactSequence}{\csname BookOriginal@trLongExactSequence\endcsname}}
\expandafter\let\csname BookOriginal@trShortExactSequence\endcsname\trShortExactSequence
\renewcommand{\trShortExactSequence}{\BookTerm{trShortExactSequence}{\csname BookOriginal@trShortExactSequence\endcsname}}
\expandafter\let\csname BookOriginal@trMorphism\endcsname\trMorphism
\renewcommand{\trMorphism}{\BookTerm{trMorphism}{\csname BookOriginal@trMorphism\endcsname}}
\expandafter\let\csname BookOriginal@trHomomorphism\endcsname\trHomomorphism
\renewcommand{\trHomomorphism}{\BookTerm{trHomomorphism}{\csname BookOriginal@trHomomorphism\endcsname}}
\expandafter\let\csname BookOriginal@trConnectingHomomorphism\endcsname\trConnectingHomomorphism
\renewcommand{\trConnectingHomomorphism}{\BookTerm{trConnectingHomomorphism}{\csname BookOriginal@trConnectingHomomorphism\endcsname}}
\expandafter\let\csname BookOriginal@trFunctor\endcsname\trFunctor
\renewcommand{\trFunctor}{\BookTerm{trFunctor}{\csname BookOriginal@trFunctor\endcsname}}
\expandafter\let\csname BookOriginal@trFunctoriality\endcsname\trFunctoriality
\renewcommand{\trFunctoriality}{\BookTerm{trFunctoriality}{\csname BookOriginal@trFunctoriality\endcsname}}
\expandafter\let\csname BookOriginal@trCategory\endcsname\trCategory
\renewcommand{\trCategory}{\BookTerm{trCategory}{\csname BookOriginal@trCategory\endcsname}}
\expandafter\let\csname BookOriginal@trCobordism\endcsname\trCobordism
\renewcommand{\trCobordism}{\BookTerm{trCobordism}{\csname BookOriginal@trCobordism\endcsname}}
\expandafter\let\csname BookOriginal@trOrientation\endcsname\trOrientation
\renewcommand{\trOrientation}{\BookTerm{trOrientation}{\csname BookOriginal@trOrientation\endcsname}}
\expandafter\let\csname BookOriginal@trCoorientation\endcsname\trCoorientation
\renewcommand{\trCoorientation}{\BookTerm{trCoorientation}{\csname BookOriginal@trCoorientation\endcsname}}
\expandafter\let\csname BookOriginal@trComplement\endcsname\trComplement
\renewcommand{\trComplement}{\BookTerm{trComplement}{\csname BookOriginal@trComplement\endcsname}}
\expandafter\let\csname BookOriginal@trModule\endcsname\trModule
\renewcommand{\trModule}{\BookTerm{trModule}{\csname BookOriginal@trModule\endcsname}}
\expandafter\let\csname BookOriginal@trFundamentalGroup\endcsname\trFundamentalGroup
\renewcommand{\trFundamentalGroup}{\BookTerm{trFundamentalGroup}{\csname BookOriginal@trFundamentalGroup\endcsname}}
\expandafter\let\csname BookOriginal@trCovering\endcsname\trCovering
\renewcommand{\trCovering}{\BookTerm{trCovering}{\csname BookOriginal@trCovering\endcsname}}
\expandafter\let\csname BookOriginal@trUniversalCover\endcsname\trUniversalCover
\renewcommand{\trUniversalCover}{\BookTerm{trUniversalCover}{\csname BookOriginal@trUniversalCover\endcsname}}
\expandafter\let\csname BookOriginal@trHomotopy\endcsname\trHomotopy
\renewcommand{\trHomotopy}{\BookTerm{trHomotopy}{\csname BookOriginal@trHomotopy\endcsname}}
\expandafter\let\csname BookOriginal@trDeformationRetract\endcsname\trDeformationRetract
\renewcommand{\trDeformationRetract}{\BookTerm{trDeformationRetract}{\csname BookOriginal@trDeformationRetract\endcsname}}
\expandafter\let\csname BookOriginal@trInterpolation\endcsname\trInterpolation
\renewcommand{\trInterpolation}{\BookTerm{trInterpolation}{\csname BookOriginal@trInterpolation\endcsname}}
\expandafter\let\csname BookOriginal@trDeformation\endcsname\trDeformation
\renewcommand{\trDeformation}{\BookTerm{trDeformation}{\csname BookOriginal@trDeformation\endcsname}}
\expandafter\let\csname BookOriginal@trMappingDegree\endcsname\trMappingDegree
\renewcommand{\trMappingDegree}{\BookTerm{trMappingDegree}{\csname BookOriginal@trMappingDegree\endcsname}}
\expandafter\let\csname BookOriginal@trIsomorphism\endcsname\trIsomorphism
\renewcommand{\trIsomorphism}{\BookTerm{trIsomorphism}{\csname BookOriginal@trIsomorphism\endcsname}}
\expandafter\let\csname BookOriginal@trEndomorphism\endcsname\trEndomorphism
\renewcommand{\trEndomorphism}{\BookTerm{trEndomorphism}{\csname BookOriginal@trEndomorphism\endcsname}}
\expandafter\let\csname BookOriginal@trAutomorphism\endcsname\trAutomorphism
\renewcommand{\trAutomorphism}{\BookTerm{trAutomorphism}{\csname BookOriginal@trAutomorphism\endcsname}}
\expandafter\let\csname BookOriginal@trGeneralLinearGroup\endcsname\trGeneralLinearGroup
\renewcommand{\trGeneralLinearGroup}{\BookTerm{trGeneralLinearGroup}{\csname BookOriginal@trGeneralLinearGroup\endcsname}}
\expandafter\let\csname BookOriginal@trStabilizer\endcsname\trStabilizer
\renewcommand{\trStabilizer}{\BookTerm{trStabilizer}{\csname BookOriginal@trStabilizer\endcsname}}
\expandafter\let\csname BookOriginal@trHomogeneousSpace\endcsname\trHomogeneousSpace
\renewcommand{\trHomogeneousSpace}{\BookTerm{trHomogeneousSpace}{\csname BookOriginal@trHomogeneousSpace\endcsname}}
\expandafter\let\csname BookOriginal@trKernel\endcsname\trKernel
\renewcommand{\trKernel}{\BookTerm{trKernel}{\csname BookOriginal@trKernel\endcsname}}
\expandafter\let\csname BookOriginal@trRank\endcsname\trRank
\renewcommand{\trRank}{\BookTerm{trRank}{\csname BookOriginal@trRank\endcsname}}
\expandafter\let\csname BookOriginal@trCokernel\endcsname\trCokernel
\renewcommand{\trCokernel}{\BookTerm{trCokernel}{\csname BookOriginal@trCokernel\endcsname}}
\expandafter\let\csname BookOriginal@trImage\endcsname\trImage
\renewcommand{\trImage}{\BookTerm{trImage}{\csname BookOriginal@trImage\endcsname}}
\expandafter\let\csname BookOriginal@trPreimage\endcsname\trPreimage
\renewcommand{\trPreimage}{\BookTerm{trPreimage}{\csname BookOriginal@trPreimage\endcsname}}
\expandafter\let\csname BookOriginal@trTensorProduct\endcsname\trTensorProduct
\renewcommand{\trTensorProduct}{\BookTerm{trTensorProduct}{\csname BookOriginal@trTensorProduct\endcsname}}
\expandafter\let\csname BookOriginal@trTorsion\endcsname\trTorsion
\renewcommand{\trTorsion}{\BookTerm{trTorsion}{\csname BookOriginal@trTorsion\endcsname}}
\expandafter\let\csname BookOriginal@trDirectSum\endcsname\trDirectSum
\renewcommand{\trDirectSum}{\BookTerm{trDirectSum}{\csname BookOriginal@trDirectSum\endcsname}}
\expandafter\let\csname BookOriginal@trLocalCoefficient\endcsname\trLocalCoefficient
\renewcommand{\trLocalCoefficient}{\BookTerm{trLocalCoefficient}{\csname BookOriginal@trLocalCoefficient\endcsname}}
\expandafter\let\csname BookOriginal@trCoherentOrientation\endcsname\trCoherentOrientation
\renewcommand{\trCoherentOrientation}{\BookTerm{trCoherentOrientation}{\csname BookOriginal@trCoherentOrientation\endcsname}}
\expandafter\let\csname BookOriginal@trSymplecticGroup\endcsname\trSymplecticGroup
\renewcommand{\trSymplecticGroup}{\BookTerm{trSymplecticGroup}{\csname BookOriginal@trSymplecticGroup\endcsname}}
\expandafter\let\csname BookOriginal@trSymplecticForm\endcsname\trSymplecticForm
\renewcommand{\trSymplecticForm}{\BookTerm{trSymplecticForm}{\csname BookOriginal@trSymplecticForm\endcsname}}
\expandafter\let\csname BookOriginal@trSymplectomorphism\endcsname\trSymplectomorphism
\renewcommand{\trSymplectomorphism}{\BookTerm{trSymplectomorphism}{\csname BookOriginal@trSymplectomorphism\endcsname}}
\expandafter\let\csname BookOriginal@trIsometry\endcsname\trIsometry
\renewcommand{\trIsometry}{\BookTerm{trIsometry}{\csname BookOriginal@trIsometry\endcsname}}
\expandafter\let\csname BookOriginal@trEigenvalue\endcsname\trEigenvalue
\renewcommand{\trEigenvalue}{\BookTerm{trEigenvalue}{\csname BookOriginal@trEigenvalue\endcsname}}
\expandafter\let\csname BookOriginal@trCharacteristicPolynomial\endcsname\trCharacteristicPolynomial
\renewcommand{\trCharacteristicPolynomial}{\BookTerm{trCharacteristicPolynomial}{\csname BookOriginal@trCharacteristicPolynomial\endcsname}}
\expandafter\let\csname BookOriginal@trMultiplicity\endcsname\trMultiplicity
\renewcommand{\trMultiplicity}{\BookTerm{trMultiplicity}{\csname BookOriginal@trMultiplicity\endcsname}}
\expandafter\let\csname BookOriginal@trEigenspace\endcsname\trEigenspace
\renewcommand{\trEigenspace}{\BookTerm{trEigenspace}{\csname BookOriginal@trEigenspace\endcsname}}
\expandafter\let\csname BookOriginal@trGeneralizedEigenspace\endcsname\trGeneralizedEigenspace
\renewcommand{\trGeneralizedEigenspace}{\BookTerm{trGeneralizedEigenspace}{\csname BookOriginal@trGeneralizedEigenspace\endcsname}}
\expandafter\let\csname BookOriginal@trComplexification\endcsname\trComplexification
\renewcommand{\trComplexification}{\BookTerm{trComplexification}{\csname BookOriginal@trComplexification\endcsname}}
\expandafter\let\csname BookOriginal@trSymmetricMatrix\endcsname\trSymmetricMatrix
\renewcommand{\trSymmetricMatrix}{\BookTerm{trSymmetricMatrix}{\csname BookOriginal@trSymmetricMatrix\endcsname}}
\expandafter\let\csname BookOriginal@trBlockDiagonalMatrix\endcsname\trBlockDiagonalMatrix
\renewcommand{\trBlockDiagonalMatrix}{\BookTerm{trBlockDiagonalMatrix}{\csname BookOriginal@trBlockDiagonalMatrix\endcsname}}
\expandafter\let\csname BookOriginal@trMaximalCompactSubgroup\endcsname\trMaximalCompactSubgroup
\renewcommand{\trMaximalCompactSubgroup}{\BookTerm{trMaximalCompactSubgroup}{\csname BookOriginal@trMaximalCompactSubgroup\endcsname}}
\expandafter\let\csname BookOriginal@trMonodromy\endcsname\trMonodromy
\renewcommand{\trMonodromy}{\BookTerm{trMonodromy}{\csname BookOriginal@trMonodromy\endcsname}}
\expandafter\let\csname BookOriginal@trForm\endcsname\trForm
\renewcommand{\trForm}{\BookTerm{trForm}{\csname BookOriginal@trForm\endcsname}}
\expandafter\let\csname BookOriginal@trVolume\endcsname\trVolume
\renewcommand{\trVolume}{\BookTerm{trVolume}{\csname BookOriginal@trVolume\endcsname}}
\expandafter\let\csname BookOriginal@trComplexStructure\endcsname\trComplexStructure
\renewcommand{\trComplexStructure}{\BookTerm{trComplexStructure}{\csname BookOriginal@trComplexStructure\endcsname}}
\expandafter\let\csname BookOriginal@trActionFunctional\endcsname\trActionFunctional
\renewcommand{\trActionFunctional}{\BookTerm{trActionFunctional}{\csname BookOriginal@trActionFunctional\endcsname}}
\expandafter\let\csname BookOriginal@trActionForm\endcsname\trActionForm
\renewcommand{\trActionForm}{\BookTerm{trActionForm}{\csname BookOriginal@trActionForm\endcsname}}
\expandafter\let\csname BookOriginal@trIntegrationCover\endcsname\trIntegrationCover
\renewcommand{\trIntegrationCover}{\BookTerm{trIntegrationCover}{\csname BookOriginal@trIntegrationCover\endcsname}}
\expandafter\let\csname BookOriginal@trInjectivityRadius\endcsname\trInjectivityRadius
\renewcommand{\trInjectivityRadius}{\BookTerm{trInjectivityRadius}{\csname BookOriginal@trInjectivityRadius\endcsname}}
\expandafter\let\csname BookOriginal@trFloerHomology\endcsname\trFloerHomology
\renewcommand{\trFloerHomology}{\BookTerm{trFloerHomology}{\csname BookOriginal@trFloerHomology\endcsname}}
\expandafter\let\csname BookOriginal@trFloerComplex\endcsname\trFloerComplex
\renewcommand{\trFloerComplex}{\BookTerm{trFloerComplex}{\csname BookOriginal@trFloerComplex\endcsname}}
\expandafter\let\csname BookOriginal@trFloerEquation\endcsname\trFloerEquation
\renewcommand{\trFloerEquation}{\BookTerm{trFloerEquation}{\csname BookOriginal@trFloerEquation\endcsname}}
\expandafter\let\csname BookOriginal@trArnoldConjecture\endcsname\trArnoldConjecture
\renewcommand{\trArnoldConjecture}{\BookTerm{trArnoldConjecture}{\csname BookOriginal@trArnoldConjecture\endcsname}}
\expandafter\let\csname BookOriginal@trMaslovIndex\endcsname\trMaslovIndex
\renewcommand{\trMaslovIndex}{\BookTerm{trMaslovIndex}{\csname BookOriginal@trMaslovIndex\endcsname}}
\expandafter\let\csname BookOriginal@trConleyZehnderIndex\endcsname\trConleyZehnderIndex
\renewcommand{\trConleyZehnderIndex}{\BookTerm{trConleyZehnderIndex}{\csname BookOriginal@trConleyZehnderIndex\endcsname}}
\expandafter\let\csname BookOriginal@trMaslovCycle\endcsname\trMaslovCycle
\renewcommand{\trMaslovCycle}{\BookTerm{trMaslovCycle}{\csname BookOriginal@trMaslovCycle\endcsname}}
\expandafter\let\csname BookOriginal@trCycle\endcsname\trCycle
\renewcommand{\trCycle}{\BookTerm{trCycle}{\csname BookOriginal@trCycle\endcsname}}
\expandafter\let\csname BookOriginal@trLoop\endcsname\trLoop
\renewcommand{\trLoop}{\BookTerm{trLoop}{\csname BookOriginal@trLoop\endcsname}}
\expandafter\let\csname BookOriginal@trLoopSpace\endcsname\trLoopSpace
\renewcommand{\trLoopSpace}{\BookTerm{trLoopSpace}{\csname BookOriginal@trLoopSpace\endcsname}}
\expandafter\let\csname BookOriginal@trFixedPoint\endcsname\trFixedPoint
\renewcommand{\trFixedPoint}{\BookTerm{trFixedPoint}{\csname BookOriginal@trFixedPoint\endcsname}}
\expandafter\let\csname BookOriginal@trCappingDisk\endcsname\trCappingDisk
\renewcommand{\trCappingDisk}{\BookTerm{trCappingDisk}{\csname BookOriginal@trCappingDisk\endcsname}}
\expandafter\let\csname BookOriginal@trCappingSphere\endcsname\trCappingSphere
\renewcommand{\trCappingSphere}{\BookTerm{trCappingSphere}{\csname BookOriginal@trCappingSphere\endcsname}}
\expandafter\let\csname BookOriginal@trCappingSurface\endcsname\trCappingSurface
\renewcommand{\trCappingSurface}{\BookTerm{trCappingSurface}{\csname BookOriginal@trCappingSurface\endcsname}}
\expandafter\let\csname BookOriginal@trChernClass\endcsname\trChernClass
\renewcommand{\trChernClass}{\BookTerm{trChernClass}{\csname BookOriginal@trChernClass\endcsname}}
\expandafter\let\csname BookOriginal@trFirstChernClass\endcsname\trFirstChernClass
\renewcommand{\trFirstChernClass}{\BookTerm{trFirstChernClass}{\csname BookOriginal@trFirstChernClass\endcsname}}
\expandafter\let\csname BookOriginal@trLagrangianGrassmannian\endcsname\trLagrangianGrassmannian
\renewcommand{\trLagrangianGrassmannian}{\BookTerm{trLagrangianGrassmannian}{\csname BookOriginal@trLagrangianGrassmannian\endcsname}}
\expandafter\let\csname BookOriginal@trGeneratingFunction\endcsname\trGeneratingFunction
\renewcommand{\trGeneratingFunction}{\BookTerm{trGeneratingFunction}{\csname BookOriginal@trGeneratingFunction\endcsname}}
\expandafter\let\csname BookOriginal@trInnerProduct\endcsname\trInnerProduct
\renewcommand{\trInnerProduct}{\BookTerm{trInnerProduct}{\csname BookOriginal@trInnerProduct\endcsname}}
\expandafter\let\csname BookOriginal@trPolarDecomposition\endcsname\trPolarDecomposition
\renewcommand{\trPolarDecomposition}{\BookTerm{trPolarDecomposition}{\csname BookOriginal@trPolarDecomposition\endcsname}}
\expandafter\let\csname BookOriginal@trSpectralFlow\endcsname\trSpectralFlow
\renewcommand{\trSpectralFlow}{\BookTerm{trSpectralFlow}{\csname BookOriginal@trSpectralFlow\endcsname}}
\expandafter\let\csname BookOriginal@trCrossing\endcsname\trCrossing
\renewcommand{\trCrossing}{\BookTerm{trCrossing}{\csname BookOriginal@trCrossing\endcsname}}
\expandafter\let\csname BookOriginal@trPseudoholomorphicCurve\endcsname\trPseudoholomorphicCurve
\renewcommand{\trPseudoholomorphicCurve}{\BookTerm{trPseudoholomorphicCurve}{\csname BookOriginal@trPseudoholomorphicCurve\endcsname}}
\expandafter\let\csname BookOriginal@trCauchyRiemann\endcsname\trCauchyRiemann
\renewcommand{\trCauchyRiemann}{\BookTerm{trCauchyRiemann}{\csname BookOriginal@trCauchyRiemann\endcsname}}
\expandafter\let\csname BookOriginal@trEnergy\endcsname\trEnergy
\renewcommand{\trEnergy}{\BookTerm{trEnergy}{\csname BookOriginal@trEnergy\endcsname}}
\expandafter\let\csname BookOriginal@trMomentum\endcsname\trMomentum
\renewcommand{\trMomentum}{\BookTerm{trMomentum}{\csname BookOriginal@trMomentum\endcsname}}
\expandafter\let\csname BookOriginal@trEnergyIdentity\endcsname\trEnergyIdentity
\renewcommand{\trEnergyIdentity}{\BookTerm{trEnergyIdentity}{\csname BookOriginal@trEnergyIdentity\endcsname}}
\expandafter\let\csname BookOriginal@trContinuationMorphism\endcsname\trContinuationMorphism
\renewcommand{\trContinuationMorphism}{\BookTerm{trContinuationMorphism}{\csname BookOriginal@trContinuationMorphism\endcsname}}
\expandafter\let\csname BookOriginal@trSimilarityPrinciple\endcsname\trSimilarityPrinciple
\renewcommand{\trSimilarityPrinciple}{\BookTerm{trSimilarityPrinciple}{\csname BookOriginal@trSimilarityPrinciple\endcsname}}
\expandafter\let\csname BookOriginal@trContinuationPrinciple\endcsname\trContinuationPrinciple
\renewcommand{\trContinuationPrinciple}{\BookTerm{trContinuationPrinciple}{\csname BookOriginal@trContinuationPrinciple\endcsname}}
\expandafter\let\csname BookOriginal@trMultiplePoint\endcsname\trMultiplePoint
\renewcommand{\trMultiplePoint}{\BookTerm{trMultiplePoint}{\csname BookOriginal@trMultiplePoint\endcsname}}
\expandafter\let\csname BookOriginal@trInfiniteOrderZero\endcsname\trInfiniteOrderZero
\renewcommand{\trInfiniteOrderZero}{\BookTerm{trInfiniteOrderZero}{\csname BookOriginal@trInfiniteOrderZero\endcsname}}
\expandafter\let\csname BookOriginal@trPerturbation\endcsname\trPerturbation
\renewcommand{\trPerturbation}{\BookTerm{trPerturbation}{\csname BookOriginal@trPerturbation\endcsname}}
\expandafter\let\csname BookOriginal@trLinearization\endcsname\trLinearization
\renewcommand{\trLinearization}{\BookTerm{trLinearization}{\csname BookOriginal@trLinearization\endcsname}}
\expandafter\let\csname BookOriginal@trBanachManifold\endcsname\trBanachManifold
\renewcommand{\trBanachManifold}{\BookTerm{trBanachManifold}{\csname BookOriginal@trBanachManifold\endcsname}}
\expandafter\let\csname BookOriginal@trFredholmOperator\endcsname\trFredholmOperator
\renewcommand{\trFredholmOperator}{\BookTerm{trFredholmOperator}{\csname BookOriginal@trFredholmOperator\endcsname}}
\expandafter\let\csname BookOriginal@trBanachSpace\endcsname\trBanachSpace
\renewcommand{\trBanachSpace}{\BookTerm{trBanachSpace}{\csname BookOriginal@trBanachSpace\endcsname}}
\expandafter\let\csname BookOriginal@trHilbertSpace\endcsname\trHilbertSpace
\renewcommand{\trHilbertSpace}{\BookTerm{trHilbertSpace}{\csname BookOriginal@trHilbertSpace\endcsname}}
\expandafter\let\csname BookOriginal@trSobolevSpace\endcsname\trSobolevSpace
\renewcommand{\trSobolevSpace}{\BookTerm{trSobolevSpace}{\csname BookOriginal@trSobolevSpace\endcsname}}
\expandafter\let\csname BookOriginal@trWeakSolution\endcsname\trWeakSolution
\renewcommand{\trWeakSolution}{\BookTerm{trWeakSolution}{\csname BookOriginal@trWeakSolution\endcsname}}
\expandafter\let\csname BookOriginal@trWeakConvergence\endcsname\trWeakConvergence
\renewcommand{\trWeakConvergence}{\BookTerm{trWeakConvergence}{\csname BookOriginal@trWeakConvergence\endcsname}}
\expandafter\let\csname BookOriginal@trConvergence\endcsname\trConvergence
\renewcommand{\trConvergence}{\BookTerm{trConvergence}{\csname BookOriginal@trConvergence\endcsname}}
\expandafter\let\csname BookOriginal@trDistribution\endcsname\trDistribution
\renewcommand{\trDistribution}{\BookTerm{trDistribution}{\csname BookOriginal@trDistribution\endcsname}}
\expandafter\let\csname BookOriginal@trEllipticRegularity\endcsname\trEllipticRegularity
\renewcommand{\trEllipticRegularity}{\BookTerm{trEllipticRegularity}{\csname BookOriginal@trEllipticRegularity\endcsname}}
\expandafter\let\csname BookOriginal@trLinearOperator\endcsname\trLinearOperator
\renewcommand{\trLinearOperator}{\BookTerm{trLinearOperator}{\csname BookOriginal@trLinearOperator\endcsname}}
\expandafter\let\csname BookOriginal@trAdjointOperator\endcsname\trAdjointOperator
\renewcommand{\trAdjointOperator}{\BookTerm{trAdjointOperator}{\csname BookOriginal@trAdjointOperator\endcsname}}
\expandafter\let\csname BookOriginal@trUnboundedOperator\endcsname\trUnboundedOperator
\renewcommand{\trUnboundedOperator}{\BookTerm{trUnboundedOperator}{\csname BookOriginal@trUnboundedOperator\endcsname}}
\expandafter\let\csname BookOriginal@trHilbertianBasis\endcsname\trHilbertianBasis
\renewcommand{\trHilbertianBasis}{\BookTerm{trHilbertianBasis}{\csname BookOriginal@trHilbertianBasis\endcsname}}
\expandafter\let\csname BookOriginal@trSemiflow\endcsname\trSemiflow
\renewcommand{\trSemiflow}{\BookTerm{trSemiflow}{\csname BookOriginal@trSemiflow\endcsname}}
\expandafter\let\csname BookOriginal@trSemigroup\endcsname\trSemigroup
\renewcommand{\trSemigroup}{\BookTerm{trSemigroup}{\csname BookOriginal@trSemigroup\endcsname}}
\expandafter\let\csname BookOriginal@trSpectrum\endcsname\trSpectrum
\renewcommand{\trSpectrum}{\BookTerm{trSpectrum}{\csname BookOriginal@trSpectrum\endcsname}}
\expandafter\let\csname BookOriginal@trUniformConvergence\endcsname\trUniformConvergence
\renewcommand{\trUniformConvergence}{\BookTerm{trUniformConvergence}{\csname BookOriginal@trUniformConvergence\endcsname}}
\expandafter\let\csname BookOriginal@trEquicontinuity\endcsname\trEquicontinuity
\renewcommand{\trEquicontinuity}{\BookTerm{trEquicontinuity}{\csname BookOriginal@trEquicontinuity\endcsname}}
\expandafter\let\csname BookOriginal@trArzelaAscoliTheorem\endcsname\trArzelaAscoliTheorem
\renewcommand{\trArzelaAscoliTheorem}{\BookTerm{trArzelaAscoliTheorem}{\csname BookOriginal@trArzelaAscoliTheorem\endcsname}}
\expandafter\let\csname BookOriginal@trImplicitFunctionTheorem\endcsname\trImplicitFunctionTheorem
\renewcommand{\trImplicitFunctionTheorem}{\BookTerm{trImplicitFunctionTheorem}{\csname BookOriginal@trImplicitFunctionTheorem\endcsname}}
\expandafter\let\csname BookOriginal@trLocalInversionTheorem\endcsname\trLocalInversionTheorem
\renewcommand{\trLocalInversionTheorem}{\BookTerm{trLocalInversionTheorem}{\csname BookOriginal@trLocalInversionTheorem\endcsname}}
\expandafter\let\csname BookOriginal@trRegularity\endcsname\trRegularity
\renewcommand{\trRegularity}{\BookTerm{trRegularity}{\csname BookOriginal@trRegularity\endcsname}}
\makeatother

